% concatenated from main.tex, sections/introduction.tex, sections/technical-overview.tex, sections/nonanytime.tex, sections/anytime.tex, sections/conclusion.tex, sections/ai-disclosure.tex, appendices/local-construction-proofs.tex, appendices/global-analysis-proofs.tex, appendices/anytime-proofs.tex
\pdfoutput=1
\documentclass[11pt,letterpaper]{article}
\usepackage[letterpaper,margin=1in]{geometry}
\usepackage[utf8]{inputenc}
\usepackage[T1]{fontenc}
\usepackage[dvipsnames]{xcolor}
\usepackage{amsmath,amssymb,amsthm,mathtools,mathrsfs}
\usepackage{enumitem,graphicx,float}
\usepackage{tikz}
\usetikzlibrary{arrows.meta,calc,positioning,decorations.pathreplacing}
\usepackage{microtype}
\microtypesetup{expansion=false}
\setlength{\emergencystretch}{2em}
\usepackage{natbib}
\setcitestyle{authoryear,open={[},close={]},citesep={;},aysep={,},yysep={,}}
\usepackage{mathptmx}
\usepackage{courier}
\usepackage{aliascnt}
\usepackage[unicode=true,colorlinks=true,breaklinks=true]{hyperref}
\definecolor{selectedpurple}{HTML}{79538A}
\definecolor{curvepurple}{HTML}{6F3F8F}
\newcommand{\yy}[1]{{\color{red}{#1}}}
\newcommand{\kz}[1]{{\color{orange}{ [\textbf{kz:}} #1]}}
\hypersetup{linkcolor=selectedpurple!85!black,citecolor=YellowOrange!85!black,
  urlcolor=Aquamarine!85!black,
  pdftitle={Silver Rate Is (Almost) Optimal for Gradient Descent Acceleration},pdfauthor={Yuhan Ye and Kaizhao Liu}}
\newtheorem{theorem}{Theorem}[section]
\newaliascnt{lemma}{theorem}
\newtheorem{lemma}[lemma]{Lemma}
\aliascntresetthe{lemma}
\newaliascnt{proposition}{theorem}
\newtheorem{proposition}[proposition]{Proposition}
\aliascntresetthe{proposition}
\newaliascnt{corollary}{theorem}
\newtheorem{corollary}[corollary]{Corollary}
\aliascntresetthe{corollary}
\newaliascnt{fact}{theorem}
\newtheorem{fact}[fact]{Fact}
\aliascntresetthe{fact}
\theoremstyle{definition}
\newaliascnt{definition}{theorem}
\newtheorem{definition}[definition]{Definition}
\aliascntresetthe{definition}
\newaliascnt{remark}{theorem}
\newtheorem{remark}[remark]{Remark}
\aliascntresetthe{remark}
\usepackage[nameinlink,capitalize]{cleveref}
\crefname{fact}{Fact}{Facts}
\Crefname{fact}{Fact}{Facts}
\crefname{appendix}{Appendix}{Appendices}
\Crefname{appendix}{Appendix}{Appendices}
\allowdisplaybreaks
\newcommand{\R}{\mathbb R}
\newcommand{\N}{\mathbb N}
\newcommand{\psil}{p_{\mathrm{sil}}}
\newcommand{\pany}{p_{\mathrm{any}}}
\DeclareMathOperator{\conv}{conv}
\DeclareMathOperator{\env}{env}
\DeclareMathOperator*{\argminop}{arg\,min}
\DeclareMathOperator*{\argmaxop}{arg\,max}

\title{Silver Rate Is (Almost) Optimal for Gradient Descent}
\author{Yuhan Ye\footnotemark[1]\\MIT\\\texttt{yyh03@mit.edu}
  \and Kaizhao Liu\thanks{Authors are listed in random order.}\\MIT\\\texttt{mrzt@mit.edu}}
\date{\today}

\begin{document}
\maketitle
\begin{abstract}
We study how far gradient descent (GD) can be accelerated by predetermined stepsizes in smooth convex optimization.
Writing $\psil=\log_2(1+\sqrt2)$, we prove an $\Omega\!\left(n^{-\psil-O(\sqrt{\log\log n/\log n})}\right)$ non-anytime lower bound.
In the anytime setting, every infinite schedule has infinitely many horizons with error $\Omega\!\left(n^{-\frac{2\psil}{1+\psil}-O(\sqrt{\log\log n/\log n})}\right)$.
Together with the silver-schedule upper bound~\citep{AltschulerParrilo2025} and the anytime upper bound~\citep{ZhangLeeDuChen2025}, our results determine the optimal polynomial convergence exponents in both settings.
\end{abstract}

\clearpage
\tableofcontents
\clearpage

\section{Introduction}
\label{sec:introduction}

We consider the unconstrained minimization of a convex smooth function $f:\R^d\to\R$.
One of the most fundamental methods for this problem is \textbf{gradient descent} (GD), which dates back to \citet{Cauchy1847}.
Given a stepsize schedule $H=(h_1,\ldots,h_n)\in[0,\infty)^n$ fixed before the optimization begins, GD iterates
\[
  x_k=x_{k-1}-h_k\nabla f(x_{k-1}),\qquad 1\le k\le n.
\]
In this paper, we study how far this basic method can be accelerated using only predetermined stepsizes.

Let $\mathcal F_L(\R^d)$ be the class of convex $L$-smooth functions on $\R^d$ with a nonempty set of minimizers.
We measure the convergence rate of a schedule $H$ by
\begin{equation}
  \mathcal R_n(H):=
  \sup_{d\in\N}\;
  \sup_{f\in\mathcal F_L(\R^d)}\;
  \sup_{x_*\in\argminop f}\;
  \sup_{x_0\in\R^d\setminus\{x_*\}}
  \frac{f(x_n)-f(x_*)}{\frac{L}{2}\lVert x_0-x_*\rVert^2}.
  \label{eq:worst-case-error}
\end{equation}
In the non-anytime (finite-horizon) setting, the schedule may be chosen separately for each fixed horizon $n$.
In the anytime setting, one infinite schedule $h=(h_k)_{k\ge1}$ is used for every horizon, with $H_n=(h_1,\ldots,h_n)$.

It is a standard textbook result that GD with the constant stepsize $1/L$ has an $O(n^{-1})$ convergence rate~\citep{LevitinPolyak1966,Nesterov2004}.
Classical acceleration methods alter the GD iteration by adding momentum or auxiliary sequences, as in the heavy-ball method of \citet{Polyak1964} and the accelerated method of \citet{Nesterov1983}.
Within the broader class of first-order methods, Nesterov's method attains the optimal $O(n^{-2})$ rate for smooth convex objectives, matching the $\Omega(n^{-2})$ oracle lower bound of \citet{NemirovskyYudin1983}.

Since GD with predetermined stepsizes is a special class of first-order method, the classical $\Omega(n^{-2})$ oracle lower bound continues to apply.
For many years, however, the standard $O(n^{-1})$ rate remained the best general upper bound known.
This leads to the fundamental question of whether a faster rate is possible.
In recent years, a growing body of work surprisingly shows that GD itself can be accelerated beyond the classical $O(n^{-1})$ rate by using carefully designed stepsize schedules that use occasional long steps and recursive structure.
The current best known non-anytime upper bound is $O(n^{-\psil})$, achieved by the silver schedule introduced by \citet{AltschulerParrilo2025}, where $\psil:=\log_2(1+\sqrt2)\approx1.2716$.
In the same work, they conjectured that $\psil$ is the optimal polynomial convergence exponent.
For a broader introduction to stepsize-based acceleration, see the recent survey~\citep{AltschulerParrilo2026Survey}.
An anytime construction based on silver-schedule blocks achieves an $O(n^{-\pany})$ rate at every stopping time, where $\pany:=2\psil/(1+\psil)\approx1.1195$~\citep{ZhangLeeDuChen2025}.
This left open whether the generic $\Omega(n^{-2})$ lower bound could be strengthened  toward these upper bounds.

Recently, a sequence of results has brought the lower bounds closer to these rates.
\citet{MaChen2026} proved an $\Omega(n^{-1.932})$ non-anytime lower bound, and a subsequent blog post by \citet{Tsai2026} sharpened it to $\Omega(n^{-\sqrt3})$.
For the anytime setting, \citet{TsaiFatkhullinZhangHe2026} proved an $\Omega(n^{-4/3})$ lower bound.\footnote{Throughout, we use an $\Omega(n^{-p})$ anytime lower bound as shorthand for the statement that no single infinite stepsize schedule achieves an $o(n^{-p})$ rate.}
In our previous work~\citep{YeLiu2026}, we proved an $\Omega(n^{-1.6342})$ non-anytime lower bound and an $\Omega(n^{-1.2408})$ anytime lower bound for arbitrary real-valued stepsizes. To our knowledge, these remain the strongest known lower bounds when negative stepsizes are allowed.
For nonnegative schedules, \citet{JungChoYun2026} further improved the non-anytime and anytime lower bounds to $\Omega(n^{-1.4500})$ and $\Omega(n^{-1.1837})$, respectively.
\subsection{Main Results}

For nonnegative schedules, we determine the optimal polynomial convergence exponents in both settings.

\begin{theorem}
\label{thm:nonanytime}
There is an absolute constant $C>0$ such that, for every sufficiently large $n$ and every schedule $H\in[0,\infty)^n$,
\begin{equation}
  \mathcal R_n(H)\ge
  n^{-\left(\psil+C\sqrt{\frac{\log\log n}{\log n}}\right)}.
  \label{eq:main-nonanytime}
\end{equation}
\end{theorem}

\begin{theorem}
\label{thm:anytime}
There is an absolute constant $C>0$ such that every infinite nonnegative schedule $H=(h_k)_{k\ge1}$ satisfies
\begin{equation}
  \mathcal R_n(H_n)\ge
  n^{-\left(\pany+C\sqrt{\frac{\log\log n}{\log n}}\right)}
  \label{eq:main-anytime}
\end{equation}
for infinitely many $n$.
\end{theorem}

Together with the corresponding upper bounds~\citep{AltschulerParrilo2025,ZhangLeeDuChen2025}, our results identify $\psil$ and $\pany$ as the optimal polynomial convergence exponents in the non-anytime and anytime settings, respectively.
This confirms the conjecture of \citet{AltschulerParrilo2025} that the silver exponent is optimal.

\begin{figure}[H]
  \centering
  \includegraphics[width=\textwidth]{figures/progress.pdf}
  \caption{Lower and upper bounds for GD with predetermined stepsizes.}
  \label{fig:progress}
\end{figure}

\subsection{Additional Related Work}
\label{sec:related-work}

As discussed above, accelerated first-order methods have a long history~\citep{Polyak1964,Nesterov1983,Nesterov2004}.
Here we focus on acceleration through time-varying stepsizes, leaving the GD update unchanged.

\paragraph{Accelerating GD via stepsize schedules.}
Stepsize-based acceleration for quadratics goes back to classical Chebyshev schedules~\citep{Young1953}.
Beyond quadratics, early work established improvements over short horizons~\citep[Chapter~8]{Altschuler2018}, followed by numerical optimization of longer schedules~\citep{DasGuptaVanParysRyu2024} and constant-factor improvements via periodic long steps under a bounded-level-set assumption~\citep{Grimmer2024}.
These analyses and searches drew on the performance estimation problem (PEP) framework~\citep{DroriTeboulle2014,TaylorHendrickxGlineur2017}.
A major breakthrough came when \citet{AltschulerParrilo2025} first established the silver rate $O(n^{-\log_2\rho})$, where $\rho=1+\sqrt{2}$, for general smooth convex objectives.
For $L$-smooth, $\mu$-strongly convex objectives, they also showed that silver stepsizes reduce the squared distance to the minimizer by a factor $\varepsilon$ in $O(\kappa^{\log_\rho 2}\log(1/\varepsilon))$ iterations, where $\kappa=L/\mu$~\citep{AltschulerParrilo2025HedgingI}.
Their work introduced silver stepsizes and recursive gluing, and conjectured the asymptotic optimality of these rates.
Subsequent work improves the constant in the function-value bound and establishes accelerated gradient-norm guarantees~\citep{GrimmerShuWang2025LongSteps}.
Further refinements of recursive gluing through composition and concatenation improved constants and optimized schedules for arbitrary prescribed horizons~\citep{GrimmerShuWang2025Composing,ZhangJiang2026}.
The silver-rate guarantee has also been extended to proximal and projected GD~\citep{BokAltschuler2026}, while random stepsizes yield asymptotic acceleration for separable smooth strongly convex objectives~\citep{AltschulerParrilo2024RandomStepsizes}.
In the anytime setting, a single predetermined infinite schedule achieves an $O(n^{-1.119})$ rate at every stopping time~\citep{ZhangLeeDuChen2025}, answering the question of whether anytime stepsize-based acceleration is possible~\citep{KornowskiShamir2024}.
For a broader overview, we refer readers to the recent survey~\citep{AltschulerParrilo2026Survey}.

\paragraph{Lower bounds.}
The classical $\Omega(n^{-2})$ first-order oracle lower bound for smooth convex optimization~\citep{NemirovskyYudin1983} applies to GD with arbitrary predetermined stepsizes.
For its standard quadratic-chain construction, see \citet[Section~2.1.2]{Nesterov2004}.
Stronger lower bounds were previously known under additional structural restrictions.
For $n=2^k-1$, \citet{WangMaYangZhou2024} determine the exact worst-case function-value bound for the original silver schedule.
Within the recursively generated class of basic $f$-composable schedules, the silver rate is also minimax optimal~\citep{GrimmerShuWang2025Composing}.
\section{Motivation and Technical Overview}
\label{sec:overview}
Replacing $f$ by $f/L$ and each $h_k$ by $Lh_k$ leaves the GD iterates and $\mathcal R_n$ unchanged, so we assume $L=1$ throughout. 
We first recall how previous constructions of \citet{MaChen2026} and \citet{JungChoYun2026} use selected steps to change the gradient direction.

\paragraph{Selecting the checkpoints.}
Fix a nonnegative schedule $h=(h_1,\ldots,h_n)$ and select indices $T=\{t_1<\cdots<t_k\}\subseteq\{1,\ldots,n\}$, which we call checkpoints.
Set $t_0=0$, $t_{k+1}=n+1$, and define
\begin{equation}
  b_i:=h_{t_i}\quad(1\le i\le k),
  \qquad
  s_i:=\sum_{t=t_{i-1}+1}^{t_i-1}h_t\quad(1\le i\le k+1).
  \label{local:checkpoint-definitions}
\end{equation}
Here $b_i$ is the selected stepsize at time $t_i$, and $s_i$ is the total stepsize of the intervening updates, indexed by $t_{i-1}<t<t_i$.
The last sum $s_{k+1}$ covers the steps after the final checkpoint.
\Cref{fig:checkpoint-gaps} illustrates this notation.

\begin{figure}[!htbp]
  \centering
  \includegraphics[width=\textwidth]{figures/checkpoint-gaps.pdf}
  \caption{The selected stepsizes are $b_i=h_{t_i}$. Each brace marks the total stepsize $s_i$ of the intervening updates.}
  \label{fig:checkpoint-gaps}
\end{figure}

For selected long steps, \citet[Theorem~4.1]{MaChen2026} construct a hard function whose gradient stays constant during the intervening updates and changes after each selected step.
See also \citet[Appendix~A]{YeLiu2026} for a geometric illustration of the trajectory.
Their hard function is globally defined depending on all the coordinates.
\citet[Section~2]{JungChoYun2026} instead construct a sum of local bivariate functions.
During the intervening updates, the current component contributes a constant gradient, while all later components remain inactive.
This gives the following improved lower bound.

\begin{fact}[{\protect\citealp{JungChoYun2026}, Lemma~2.3}]
\label{fact:jcy-transfer}
For every positive schedule $h$ and every checkpoint set $T$,
\begin{equation}
  \mathcal R_n(h)\ge
  \frac{1}{4(1+s_{k+1})}
  \prod_{i=1}^k
  \left[\frac{b_i}{2(2+s_i)}\right]^2.
  \label{eq:jcy-transfer}
\end{equation}
For $k=0$, the product is empty and $s_1=\sum_{t=1}^n h_t$.
\end{fact}

Specifically, their construction uses the one-sided Huber function
\begin{equation*}
  H_\delta(z):=
  \begin{cases}
    0,&z\le0,\\
    z^2/2,&0\le z\le\delta,\\
    \delta z-\delta^2/2,&z\ge\delta,
  \end{cases}
  \qquad \delta>0,
\end{equation*}
and joins adjacent coordinates through
\begin{equation*}
  F(x):=\frac14\sum_{i=1}^k
  H_{\delta_i}\bigl(x^{(i)}-x^{(i+1)}-\tau_i\bigr)
  +\frac12H_{\delta_{k+1}}\bigl(x^{(k+1)}-\tau_{k+1}\bigr).
\end{equation*}
See~\citep[Equation (4)]{JungChoYun2026}.
The parameters $\delta_i>0$ and activation thresholds $\tau_i\ge0$ are chosen based on the stepsizes and the selected checkpoints~\citep[Equation (27)]{JungChoYun2026}.
Starting from the first coordinate, the selected step of size $b_i$ activates component $i+1$, while all later components remain inactive.\footnote{A component is active at an iterate if its gradient there is nonzero, and inactive otherwise.}
This step raises coordinate $i+1$ above its activation threshold, allowing the next component to contribute to the gradient.
The final one-dimensional component gives the function-value lower bound in \eqref{eq:jcy-transfer}.
They then optimize over checkpoint sets and analyze the resulting bound recursively, obtaining the $\Omega(n^{-\log_2(1+\sqrt3)})$ non-anytime lower bound~\citep[Section~3]{JungChoYun2026}.

\paragraph{The bottleneck.}
The coefficient $2$ of $s_i$ in \eqref{eq:jcy-transfer} suggests where to improve the construction.
During the intervening updates, component $i$ stays in its affine region and contributes a gradient proportional to $e_i-e_{i+1}$.
Its contribution decreases coordinate $i$ and increases coordinate $i+1$ by the same amount.
Both motions decrease the difference $x^{(i)}-x^{(i+1)}$ that keeps the component active.
To keep component $i+1$ inactive during the intervening updates, its activation threshold is set to the value of coordinate $i+1$ just before the update with stepsize $b_i$.
Only the additional increase produced by this update is then passed to the next component.
The resulting transfer ratio is $b_i/[c(2+s_i)]$, with $c=2$.
Following the exponent calculation in \citet[Section~3]{JungChoYun2026}, a coefficient $c$ of $s_i$ suggests a lower bound of order $n^{-\log_2(1+\sqrt{1+c})}$.
This motivates our two-coordinate hard function, which changes the gradient direction during the intervening updates to make $c$ arbitrarily close to $1$, improving the lower-bound exponent toward the silver exponent.

\paragraph{Our innovation: switching the gradient direction during the intervening updates.}
To reduce the coefficient $c$ of $s_i$, we modify the local hard function on coordinates $i$ and $i+1$.
We keep the initial gradient horizontal and then gradually turn it toward the negative $x^{(i+1)}$ direction.\footnote{Here horizontal means that the gradient has zero component in coordinate $i+1$, while vertical means that it has zero component in coordinate $i$.}

Let $\epsilon,\gamma>0$ be parameters to be chosen later, and set
\[
  c=\sqrt{1+2\gamma+\epsilon^2},\qquad
  R=\sqrt{\epsilon^2+(1+\gamma)^2}.
\]
We choose gradient vectors on the circular arc of radius $R$ centered at $(0,\gamma)$, joining $(c,0)$ to $(\epsilon,-1)$ in the quadrant $\{v^{(1)}\ge0,v^{(2)}\le0\}$.
Here $v=(v^{(1)},v^{(2)})$ is a two-coordinate gradient vector.
Its two coordinates will correspond to coordinates $i$ and $i+1$ of the full construction.
Near time $t_i$, the gradient turns toward $(\epsilon,-1)$.
For small $\epsilon$, the selected step of size $b_i$ therefore increases coordinate $i+1$ much more than it decreases coordinate $i$.
Let $K$ be the convex hull of this arc and the origin.

To obtain a GD trajectory for the prescribed stepsizes, we construct it backward.
Write $\alpha_1,\ldots,\alpha_m$ for the steps in this two-coordinate problem, $r_j$ for its iterate after $j$ steps, and $v_j=\Pi_K(r_j)$ for the corresponding gradient, where $\Pi_K$ denotes Euclidean projection onto $K$.
We fix the point and gradient just before the selected update to be $r_m=v_m=(\epsilon,-1)$.
For $j=m,m-1,\ldots,1$, choose the preceding point and gradient to satisfy
\[
  r_{j-1}=r_j+\alpha_jv_{j-1},
  \qquad v_{j-1}=\Pi_K(r_{j-1}).
\]
The first identity is the GD update read backward.
Once the backward gradient reaches $(c,0)$, all earlier gradients remain horizontal.
If the intervening updates have small total stepsize, the first gradient may lie partway along the arc. The total stepsize spent turning is can be controlled. This is made precise in \Cref{lem:local-backward}.

For each pair of adjacent coordinates $x^{(i)}$ and $x^{(i+1)}$, we use the vectors constructed above to form a finite maximum of affine functions.
Its Moreau envelope is convex and $1$-smooth and has the prescribed gradients along the constructed trajectory. This is made precise in \Cref{lem:local-transfer}.
We then scale and translate the resulting function to obtain $\Phi_i$ as defined in \eqref{local:chain-component}.
The details are given in \cref{sec:local}.
\Cref{fig:gradient-turning} shows the local gradients $\nabla\Phi_i(x_t)$, up to a common positive scale, on the left and the resulting iterate path on the right.
The gray region marks where this component vanishes.
The next component has activation level $\ell_{i+1}=x_{t_i-1}^{(i+1)}$, the value of coordinate $i+1$ just before taking the step of size $b_i$.
From $x_{t_i}$ onward, that coordinate may decrease as the next component runs, but it never falls below $\ell_{i+1}$. The old component therefore remains inactive.

\begin{figure}[!htbp]
  \centering
  \includegraphics[width=\textwidth]{figures/gradient-turning.pdf}
  \caption{Local gradients $\nabla\Phi_i(x_t)$ and iterates. The step $h_{t_i}$ (orange) activates the next component.}
  \label{fig:gradient-turning}
\end{figure}

\begin{samepage}
\paragraph{The improved hard-function bound.}
Let $\ell_i$ be the activation level of component $i$.
At $x_{t_{i-1}}$, $D_i=x_{t_{i-1}}^{(i)}-\ell_i$ is the amount by which coordinate $i$ exceeds this level.
We start with $D_1=1$.
After scaling and joining the components, the step of size $b_i$ leaves coordinate $i+1$ a distance $D_{i+1}=D_iG_i$ above $\ell_{i+1}$, where
\[
  G_i=\frac{b_i}{c(a+s_i)+\epsilon b_i}.
\]
The constant $a$ depends only on $\epsilon,\gamma$.
The coefficient of $s_i$ is now $c$, which can be made arbitrarily close to $1$, instead of $2$.
The final bound introduce a constant $a$, a lower bound on $b_i$, and the additional term $\epsilon b_i$, made precise in the following lemma.
\par
\end{samepage}

\begin{lemma}[Checkpoint transfer bound]
\label{thm:transfer}
Fix $\epsilon,\gamma>0$ and let $c:=\sqrt{1+2\gamma+\epsilon^2}$.
There are finite constants $a=a(\epsilon,\gamma)$ and $B_0=B_0(\epsilon)$ such that, for every nonnegative schedule and every checkpoint set whose selected steps satisfy $b_i\ge B_0$,
\begin{equation}
  \mathcal R_n(h)\ge
  \frac{1}{4(1+s_{k+1})}
  \prod_{i=1}^k
  \left[\frac{b_i}{c(a+s_i)+\epsilon b_i}\right]^2.
  \label{eq:transfer}
\end{equation}
For $k=0$, the product is empty and $s_1=\sum_{t=1}^n h_t$.
\end{lemma}

We explain the construction in \cref{sec:local} and give the detailed verifications in \cref{app:local}.
The constants $a$ and $B_0$ may grow as $\epsilon,\gamma$ decrease, so we must keep track of them when choosing the parameters.

\paragraph{From the construction to the main lower bounds.}
To use this bound, we must choose checkpoints that keep the product $\prod_{i=1}^kG_i$ large relative to the total stepsize $s_{k+1}$.
We express this choice as minimizing \eqref{global:eq-checkpoint-expression} over checkpoint sets and bound its minimum recursively (\cref{global:lem-value-decomposition,thm:weighted-recursion}), following the idea of \citet[Section~3]{JungChoYun2026}.
The main difference is the additional term $\epsilon b_i$ in \eqref{eq:transfer}.
Although $\epsilon$ is small, $b_i$ has no upper bound, so $\epsilon b_i$ need not be small relative to $c(a+s_i)$ and cannot be absorbed into the constant $a$ independently of the schedule.
In the recursion, these terms shift the parameters of the shorter problems, as shown in \eqref{global:eq-value-decomposition}.
To handle the additional $\epsilon b_i$ terms, we use the splitting inequality in \cref{global:lem-gap-splitting}, which retains the correction term $\sum_i(\epsilon b_i/c)^\nu$ for a suitable exponent $\nu\in[1/2,1)$.
This term cancels the extra terms arising from the shifted parameters in the induction, yielding the bound in \cref{thm:weighted-recursion}.
This gives an $\Omega(n^{-p})$ non-anytime lower bound for every $p>\psil$.
Keeping track of the constants as $p\downarrow\psil$ gives \cref{thm:nonanytime}, as shown in \cref{sec:global}.

For the anytime result, we follow the approaches of \citet[Section~4.1]{TsaiFatkhullinZhangHe2026} and \citet[Section~4]{JungChoYun2026}.
At horizons where the last step is the largest so far, we compare the total stepsize with that largest step and the final error.
Using our stronger construction gives the exponent $2p/(1+p)$ at infinitely many horizons.
Tracking the constants as $p\downarrow\psil$ then yields \cref{thm:anytime}, as shown in \cref{sec:anytime}.

\section{(Almost) Tight Non-Anytime Lower Bound}
\label{sec:nonanytime}

We first explain how to construct the hard functions in \cref{thm:transfer}, and then give the lemmas that complete the proof of \cref{thm:nonanytime}.
The detailed proofs and verifications are in \cref{app:local} and \cref{app:global}, respectively.

\subsection{Proof of Lemma~\ref{thm:transfer}}
\label{sec:local}

Fix $\epsilon,\gamma>0$ and let $c=\sqrt{1+2\gamma+\epsilon^2}$.
The construction has two parts.
First, we build a two-coordinate function for the intervening updates and the selected step that follows them.
We then join copies of this function, keeping earlier and later components inactive while the current one runs.

\paragraph{Choosing the gradient directions.}
Use the arc and convex set $K$ from \cref{sec:overview}.
Because $\nabla\operatorname{env}_1\sigma_K(r)=\Pi_K(r)$, the backward construction must satisfy both the GD update and the projection condition.
The following lemma guarantees this for any given stepsizes $\alpha_j$ and bounds the total displacement of the coordinates.

\begin{lemma}
\label{lem:local-backward}
There is a finite constant $H=H(\epsilon,\gamma)$ such that, for every $m\ge0$ and every sequence of stepsizes $\alpha_1,\ldots,\alpha_m\ge0$, there are points $r_0,\ldots,r_m$ and vectors $v_0,\ldots,v_m$ on the arc satisfying
\[
  r_m=v_m=(\epsilon,-1),\qquad
  r_j=r_{j-1}-\alpha_jv_{j-1}\quad(1\le j\le m),
  \qquad v_j=\Pi_K(r_j)\quad(0\le j\le m).
\]
Moreover, define
\begin{equation}
  S:=\sum_{j=1}^m\alpha_j,\qquad
  P:=\sum_{j=1}^m\alpha_j v_{j-1}^{(1)},\qquad
  Q:=-\sum_{j=1}^m\alpha_jv_{j-1}^{(2)},
  \label{local:displacement-constants}
\end{equation}
then the construction also satisfies
\begin{equation}
P\le cS,\qquad
    0\le Q \le H,
  \label{local:displacement-bounds}
\end{equation}
\end{lemma}

Here, the shift $\gamma>0$ allows the gradient direction to change as we construct the trajectory backward.
If $\gamma=0$, the gradient would remain at $(\epsilon,-1)$.
Keeping $\epsilon>0$ lets a sufficiently large $b_i$ move the first coordinate far enough to put the old component into a region where it remains zero as the next component runs.
The proof and explicit formulas for $H$ are given in \cref{local:app-backward}.

\paragraph{Constructing the local hard function.}
\Cref{lem:local-backward} gives the trajectory during the intervening updates.
We now scale and translate it to start at $(1,0)$, then perform a final GD update with stepsize $b$.
The final update must raise the second coordinate by a specified amount and leave the old component zero while the next component runs.
The following lemma records both properties.

\begin{lemma}[Two-coordinate transfer]
\label{lem:local-transfer}
There is a finite constant $C_0=C_0(\epsilon,\gamma)$ with the following property.
Let $m\ge0$, $\alpha_1,\ldots,\alpha_m\ge0$, $S=\sum_{j=1}^m\alpha_j$, $b\ge1+\epsilon^{-2}$, and $A\ge C_0$.
There are a nonnegative, convex, $1$-smooth function $f:\R^2\to\R$ and numbers $U,Y\ge0$ such that GD follows the trajectory below, where $G=b/(cS+\epsilon b+A)$:
\[
  (1,0)\xrightarrow{\alpha_1,\ldots,\alpha_m}(U+\epsilon G,Y)
  \xrightarrow{b}(U,Y+G).
\]
During the first $m$ updates, the first coordinate is nonincreasing and the second is nondecreasing.
Moreover, the function satisfies
\begin{equation}
  f(x,y)=0\quad\text{if }x\le0,\ y\ge0,
  \qquad
  f(U,y)=0\quad\text{if }y\ge Y.
  \label{local:zero-regions}
\end{equation}
\end{lemma}
In the construction, the displacement bounds \eqref{local:displacement-bounds}, together with $b\ge1+\epsilon^{-2}$ and $A\ge C_0$, ensure \eqref{local:zero-regions}; see \cref{local:app-transfer}.
The proof and explicit formulas for $C_0$ are given in \cref{local:app-transfer}.

We are going to pass the additional height $G$ to the next component.
When we join the components, the first zero region keeps a component inactive before its turn, and the second keeps it inactive afterward, provided that its first coordinate stays at $U$ and its second coordinate stays at least $Y$.
In the next paragraph, we will see that the monotonicity of the first $m$ updates ensures the above condition.


\paragraph{Combining the local hard functions.}
For each $i$, we apply \cref{lem:local-transfer} with the local stepsizes specified below to obtain a function $f_i$.
We scale and translate $f_i$ to form $\Phi_i$ in \eqref{local:chain-component}, then define the hard function $F$ as half the sum of these components and a final one-dimensional Huber function in \eqref{local:global-objective}.
Each coordinate appears in at most two $1$-smooth terms, so this makes the full objective $F$ $1$-smooth.
Because of the factor $1/2$, a GD step of size $h_t$ corresponds to a local step of size $h_t/2$.

Choose $B_0=2(1+\epsilon^{-2})$ and $a\ge 2C_0/c$.
For component $i$, put $m=t_i-t_{i-1}-1$ and apply \cref{lem:local-transfer} with
\[
  \alpha_j=\frac{h_{t_{i-1}+j}}2\quad(1\le j\le m),\qquad
  S=\frac{s_i}2,\qquad b=\frac{b_i}2,\qquad A=\frac{ca}2.
\]
Its transfer ratio is
\[
  G_i=\frac{b_i}{c(a+s_i)+\epsilon b_i}.
\]
Write $f_i,U_i,Y_i$ for the corresponding function and coordinates supplied by \cref{lem:local-transfer}. Set
\begin{equation}
  D_1:=1,\quad\ell_1:=0,\qquad
  D_{i+1}:=D_iG_i,\quad\ell_{i+1}:=D_iY_i\quad(1\le i\le k).
  \label{local:chain-parameters}
\end{equation}
Here, $D_i$ is the amount by which coordinate $i$ exceeds its activation level $\ell_i$ at $x_{t_{i-1}}$.
For $x\in\R^{k+1}$, define the component acting on coordinates $i,i+1$ by
\begin{equation}
  \Phi_i(x):=D_i^2 f_i\left(
    \frac{x^{(i)}-\ell_i}{D_i},\frac{x^{(i+1)}}{D_i}\right).
  \label{local:chain-component}
\end{equation}
This scaling preserves $1$-smoothness, while the first input measures height above $\ell_i$.
The final objective is defined as half the sum of these components and a final one-dimensional Huber function:
\begin{equation}
  F(x):=\frac12\sum_{i=1}^k\Phi_i(x)
       +\frac12H_\delta(x^{(k+1)}-\ell_{k+1}).
  \label{local:global-objective}
\end{equation}
From $x_{t_i}$ onward, the subsequent iterates satisfy
\[
  x^{(i)}=\ell_i+D_iU_i,\qquad
  x^{(i+1)}\ge\ell_{i+1}=D_iY_i.
\]
Thus the two inputs to $f_i$ remain $U_i$ and at least $Y_i$, respectively, so the second zero region in \Cref{lem:local-transfer} keeps $\Phi_i$ inactive.
The first zero region keeps later components inactive until their turn.
The induction verifying these coordinate relations is in \cref{local:app-chain}.

At $x_{t_k}$, the last coordinate is a distance $D=\prod_{i=1}^kG_i$ above $\ell_{k+1}$.
With $\delta=D/(1+s_{k+1})$, the final Huber calculation gives
\[
  F(x_n)-\min F
  =\frac{D^2}{4(1+s_{k+1})}
  =\frac{1}{4(1+s_{k+1})}
    \prod_{i=1}^k
    \left[\frac{b_i}{c(a+s_i)+\epsilon b_i}\right]^2.
\]
The construction starts at $e_1$ and has a minimizer at the origin, so the initial distance is $1$.
Under the normalization in \eqref{eq:worst-case-error}, $\mathcal R_n(h)$ is at least twice this value.
Dropping the factor $2$ gives \cref{thm:transfer}.
When there are no checkpoints, the same Huber calculation applies with $D=1$.
The smoothness of $F$ and final-value calculations are also verified in \cref{local:app-chain}.

\begin{samepage}
\subsection{Proof of Theorem~\ref{thm:nonanytime}}
\label{sec:global}
\label{sec:silver}

In this section, we sketch the proof of \cref{thm:nonanytime}.
The main idea of our proof follows \citet[Section~3 and Appendix~C]{JungChoYun2026}.
In our lower bound \eqref{eq:transfer}, the denominator $c(a+s_i)+\epsilon b_i$ replaces their $2(2+s_i)$.
The parameter choice in \eqref{global:eq-parameter-choice} allows $c>1$ to be arbitrarily close to $1$, and the resulting smaller coefficient of $s_i$ yields a stronger lower bound approaching the silver rate.
To handle the additional $\epsilon b_i$ terms, we retain a correction term in the splitting inequality of \cref{global:lem-gap-splitting}. This is the main difference between our proof technique and the one used in \citet[Section~3]{JungChoYun2026}.
The full details are given in \cref{app:global}.
\par
\end{samepage}

Throughout this subsection, fix $0<\epsilon<1/2$ and $\gamma>0$, set $c:=\sqrt{1+2\gamma+\epsilon^2}$, and choose the constants in \cref{thm:transfer} with $a\ge2B_0+2$ and $a\ge2$.

For a fixed schedule $h\in[0,\infty)^n$, we optimize the lower bound \eqref{eq:transfer} over the checkpoints and then seek a bound uniform over all schedules.
Using the notation $T,b_i,s_i$ from \eqref{local:checkpoint-definitions}, define
\begin{equation}
  \Psi_\lambda(T;h):=(\lambda+s_{k+1})
  \prod_{i=1}^k\left(\epsilon+\frac{c(a+s_i)}{b_i}\right),
  \label{global:eq-checkpoint-expression}
\end{equation}
\begin{equation*}
  V_\lambda(h):=\min_{T\subseteq\{1,\ldots,n\}}\Psi_\lambda(T;h),
  \qquad
  U_n(\lambda):=\sup_{h\in[0,\infty)^n}V_\lambda(h).
\end{equation*}
Here $\lambda\ge a$ is a hyperparameter, chosen later.
For the empty set, $\Psi_\lambda(\varnothing;h)=\lambda+\sum_{t=1}^n h_t$, and $U_0(\lambda)=\lambda$.\footnote{We set $\Psi_\lambda(T;h)=+\infty$ if $T$ selects a zero step, and interpret a minimum over no indices as $+\infty$.}
The calculation in \cref{lem:cost-lower-bound} gives
\begin{equation}
  \mathcal R_n(h)\ge\frac{\lambda-1}{V_\lambda(h)^2}
  \ge\frac{\lambda-1}{U_n(\lambda)^2}.
  \label{global:eq-cost-lower-bound}
\end{equation}
Thus, minimizing $\Psi_\lambda$ gives a lower bound for a fixed schedule, and bounding $U_n(\lambda)$ makes it uniform over schedules.
The following identity reduces the problem to shorter sequences.

\begin{lemma}[Exact recursive decomposition]
\label{global:lem-value-decomposition}
For every $n\ge1$, $h\in[0,\infty)^n$, and $\lambda\ge a$,
\begin{equation}
  V_\lambda(h)
  =\min\left\{
    \lambda+\sum_{j=1}^n h_j,\;
    \min_{\substack{1\le t\le n\\ h_t>0}}
    \frac{c}{h_t}
    V_{a+\epsilon h_t/c}(h_{1:t-1})V_\lambda(h_{t+1:n})
  \right\}.\footnote{Compared with \citet[Lemma~3.1]{JungChoYun2026}, the coefficient $2$ becomes $c$, and the left parameter $2$ becomes $a+\epsilon h_t/c$.
The proof uses the same factorization.}
  \label{global:eq-value-decomposition}
\end{equation}
Here $h_{i:j}:=(h_i,\ldots,h_j)$ and $V_\mu(\varnothing)=\mu$.
\end{lemma}

\paragraph{Proof sketch of \cref{global:lem-value-decomposition}.}
\pushQED{\qed}
Splitting a nonempty checkpoint set at $t\in T$ with $h_t>0$ gives
\[
  \Psi_\lambda(T;h)
  =\frac{c}{h_t}\,
    \Psi_{a+\epsilon h_t/c}(T_L;h_{1:t-1})
    \Psi_\lambda(T_R;h_{t+1:n}),
\]
where $T_L$ and $T_R$ are the checkpoints before and after $t$, reindexed within each subsequence.
Minimize independently over the two sets, then over $t$, and include the empty set.
The factor calculation is expanded in \cref{global:app-value-decomposition}.
\popQED\\

We next bound $U_n(\lambda)$.
In \eqref{global:eq-value-decomposition}, the additional $\epsilon$ term shifts the parameter of the subproblem on $h_{1:t-1}$ from $a$ to $a+\epsilon h_t/c$.
To account for it, fix $\nu\in[1/2,1)$ and set
\begin{equation}
  \kappa:=\left(\frac{\epsilon}{c-\epsilon}\right)^\nu.
  \label{global:eq-kappa-choice}
\end{equation}
Consider the following scalar condition: whenever $x,y\ge a$ and $B>0$ satisfy $w=x+y+B-a$ and $wB=(c-\epsilon)xy$, we have
\begin{equation}
  w^\nu+\kappa B^\nu\le x^\nu+y^\nu.
  \label{global:eq-scalar-hypothesis}
\end{equation}
Under this condition, we obtain the following bound on $U_n(\lambda)$.\footnote{\cref{thm:weighted-recursion} has the same form as \citet[Lemma~3.2]{JungChoYun2026}.
Here the scalar hypothesis has the additional term $\kappa B^\nu$.}

\begin{lemma}
\label{thm:weighted-recursion}
Let $\nu\in[1/2,1)$.
If the scalar hypothesis \eqref{global:eq-scalar-hypothesis} holds with \eqref{global:eq-kappa-choice}, then, for every $n\ge0$ and $\lambda\ge a$,
\begin{equation}
  U_n(\lambda)^\nu\le a^\nu n+\lambda^\nu.
  \label{eq:cost-growth}
\end{equation}
\end{lemma}

\paragraph{Proof sketch of \cref{thm:weighted-recursion}.}
\pushQED{\qed}
First, for $n\ge1$, choose a schedule $h$ attaining $U_n(\lambda)$ as guaranteed by \cref{global:lem-extremal-schedule}.
The same lemma gives a nonempty set $T=\{t_1<\cdots<t_k\}$ such that both $T$ and the empty set minimize $\Psi_\lambda$.
Evaluating $\Psi_\lambda$ at these two sets gives $U_n(\lambda)=\lambda+\sum_t h_t=\Psi_\lambda(T;h)$.

The splitting inequality in \cref{global:lem-gap-splitting} adapts the perturbation argument of \citet[Appendix~C.6]{JungChoYun2026} to retain the correction term $\sum_i(\epsilon b_i/c)^\nu$.
This is where \eqref{global:eq-scalar-hypothesis} is needed.
Applying that lemma gives
\begin{equation}
  U_n(\lambda)^\nu+\sum_{i=1}^k(\epsilon b_i/c)^\nu
  \le(\lambda+s_{k+1})^\nu+\sum_{i=1}^k(a+s_i+\epsilon b_i/c)^\nu.
  \label{global:eq-gap-splitting}
\end{equation}
The identities \eqref{global:eq-gap-identities} in \cref{global:lem-gap-splitting} show that $a+s_i+\epsilon b_i/c$ equals $V_{a+\epsilon b_i/c}$ evaluated on the intervening stepsizes between checkpoints $i-1$ and $i$.
Similarly, $\lambda+s_{k+1}$ equals $V_\lambda$ evaluated on the stepsizes after the last checkpoint.

We now induct on $n$, simultaneously for all $\lambda\ge a$, starting from $U_0(\lambda)=\lambda$.
Each intervening sequence has fewer than $n$ steps, and together they contain $n-k$ steps.
By definition, $V_\mu$ evaluated on a sequence of $m$ stepsizes is at most $U_m(\mu)$.
Applying the induction hypothesis gives
\[
\begin{aligned}
  U_n(\lambda)^\nu+\sum_{i=1}^k(\epsilon b_i/c)^\nu
  &\le a^\nu(n-k)+\lambda^\nu+\sum_{i=1}^k(a+\epsilon b_i/c)^\nu\\
  &\le a^\nu n+\lambda^\nu+\sum_{i=1}^k(\epsilon b_i/c)^\nu,
\end{aligned}
\]
where the last inequality uses $(a+\epsilon b_i/c)^\nu\le a^\nu+(\epsilon b_i/c)^\nu$.
Cancelling the identical sums proves \eqref{eq:cost-growth}.
\popQED\\

It remains to verify \eqref{global:eq-scalar-hypothesis} for $p=1/\nu$ close to $\psil$.
In the limiting case $c=1$, $\epsilon=0$, the silver inequality
\begin{equation}
  z^{\psil}+(1-z)^{\psil}+[z(1-z)]^{\psil}\le1,
  \qquad z\in[0,1],
  \label{global:eq-silver-inequality}
\end{equation}
gives this condition at $p=\psil$, as shown in \cref{global:app-scalar}.
Writing $\rho=1+\sqrt2$, the identity $\rho^2=2\rho+1$ makes the inequality tight at the balanced split $z=1/2$.
Its proof follows the concavity argument of \citet[Lemma~C.4]{JungChoYun2026}, with coefficient $1$ instead of $2$.
For $\epsilon>0$, the next lemma quantifies the resulting increase in the exponent.
Write $t_+=\max\{t,0\}$.

\begin{lemma}
\label{global:lem-scalar-condition}
Let \(p\in[\psil,3/2]\), \(\nu=1/p\), \(0<c-\epsilon\le2\), and \(0<\epsilon<1/2\).
Suppose that \(\tau:=\epsilon^{1/p}\le1/10\) and
\begin{equation}
\label{global:eq-scalar-sufficient}
 p-\psil\ge(c-\epsilon-1)_++6\epsilon^{1/p}.
\end{equation}
Then the scalar hypothesis \eqref{global:eq-scalar-hypothesis} holds.
\end{lemma}

The proof is provided in \cref{global:app-scalar}.

\paragraph{Proof of Theorem~\ref{thm:nonanytime}.}
\pushQED{\qed}
Fix $0<\eta\le1/10$ and choose
\begin{equation}
\label{global:eq-parameter-choice}
 \begin{gathered}
 p:=\psil+\eta,\qquad \gamma:=\frac{\eta}{4},\qquad \epsilon:=\left(\frac{\eta}{16}\right)^p,\\
 c:=\sqrt{1+2\gamma+\epsilon^2}.
 \end{gathered}
\end{equation}
These parameters satisfy \cref{global:lem-scalar-condition}, as verified in \cref{global:app-parameters}.

Set $\nu=1/p$ and, for $n\ge1$, choose $\lambda=an^p$.
By \cref{thm:weighted-recursion}, $V_\lambda(h)\le U_n(\lambda)\le2^pan^p$.
Since $a\ge2$, we have $\lambda-1\ge an^p/2$, so \eqref{global:eq-cost-lower-bound} yields
\begin{equation}
  \mathcal R_n(h)\ge\frac{1}{2^{2p+1}a}\,n^{-p}
  \qquad\text{for every }h\in[0,\infty)^n.
  \label{eq:fixed-p-rate}
\end{equation}
The bounds on the construction constants in \cref{global:app-parameters} give
\begin{equation}
\label{eq:parameter-growth}
 \log a=O\!\left(\eta^{-1}\log(1/\eta)\right), \qquad 0<\eta\le1/10,
\end{equation}
with an absolute implied constant.
Combining this estimate with \eqref{eq:fixed-p-rate} yields
\begin{equation*}
 \mathcal R_n(h)\ge n^{-\psil} \exp\!\left\{-\eta\log n-O\!\left(\eta^{-1}\log(1/\eta)\right)\right\}.
\end{equation*}
The implied constant is absolute and the estimate is uniform over $h$.
For sufficiently large $n$, take
\begin{equation*}
 \eta:=\sqrt{\frac{\log\log n}{\log n}}\le\frac{1}{10}.
\end{equation*}
Both terms in the exponent are then $O(\sqrt{\log n\,\log\log n})$.
Thus, for an absolute constant $C>0$ and every schedule $h\in[0,\infty)^n$,
\begin{equation*}
 \mathcal R_n(h)\ge n^{-\psil} \exp\!\left\{-C\sqrt{\log n\,\log\log n}\right\} =n^{-\left(\psil+C\sqrt{\frac{\log\log n}{\log n}}\right)}.
\end{equation*}
This proves \cref{thm:nonanytime}.
\popQED

\section{(Almost) Tight Anytime Lower Bound}
\label{sec:anytime}

We now use the same hard functions to prove \cref{thm:anytime}.
Following the record-time approach of \citet[Section~4.1]{TsaiFatkhullinZhangHe2026}, also used by \citet[Section~4]{JungChoYun2026}, we compare the total stepsize with the largest step seen so far.

Fix $p\in(\psil,\psil+1/10]$, set $\nu:=1/p$, and use the parameters $a,c,\epsilon,B_0$ from \cref{global:app-parameters}, with $\eta=p-\psil$.
For each prefix $H_n=(h_1,\ldots,h_n)$, write
\[
    S_n:=\sum_{t=1}^n h_t,
    \qquad M_n:=\max_{1\le t\le n}h_t,
    \qquad r_n:=\mathcal R_n(H_n).
\]
We call $n$ a \emph{record time} if $h_n=M_n$.
The following estimate is the main step.

\begin{lemma}
\label{any:lem:record-sum}
For every $p\in(\psil,\psil+1/10]$, there are constants $C_p<\infty$ and $r_{0,p}>0$ such that every record time $n$ with $M_n\ge B_0$ and $r_n\le r_{0,p}$ satisfies
\begin{equation}
    S_n\le C_p nM_n^{1-\nu}.
    \label{any:eq:record-sum}
\end{equation}
\end{lemma}

The proof of \cref{any:lem:record-sum} is given in \cref{any:app-record-sum}. The full proof of \cref{thm:anytime} is in \cref{any:app-main-proof}.

\section{Concluding Remarks}
\label{sec:conclusion}

In this paper, we proved an almost tight lower bound for GD with predetermined nonnegative stepsizes.
Our key innovation is to bend the gradient direction during the local updates. 
Our lower bounds match the upper bounds in both non-anytime and anytime cases, up to subpolynomial terms.

The remaining open questions are:

\begin{enumerate}
\item Can the lower and upper bounds be matched for schedules that may include negative stepsizes?
To our knowledge, the best known lower bounds for such schedules are $\Omega(n^{-1.6342})$ in the non-anytime case and $\Omega(n^{-1.2408})$ in the anytime case, as established in our earlier paper~\citep{YeLiu2026}.
Extending the lower bounds proved here to such schedules remains open, although we believe that our approach can be adapted to handle negative stepsizes.

\item Can the $\sqrt{\log\log n/\log n}$ losses in the exponents of \eqref{eq:main-nonanytime} and \eqref{eq:main-anytime} be reduced or removed?\footnote{We have obtained preliminary improvements by modifying our existing construction to choose gradient vectors along a logarithmic spiral rather than the circular arc in \eqref{local:arc}.}
\end{enumerate}

\section*{Acknowledgement}
We thank Jason Altschuler and Pablo Parrilo for helpful feedback regarding the literature on stepsize-based acceleration for GD.
\section*{AI Disclosure}

We carefully read the paper by \citet{JungChoYun2026} and traced the difference between their exponent $\log_2(1+\sqrt3)$ and the silver exponent $\log_2(1+\sqrt2)$ to the factor $2$ in their local transfer estimate.
Their local Huber component has a fixed gradient direction.
It changes both adjacent coordinates at the same speed, limiting the room above the activation level for the next component.
We wondered whether we could increase this room by bending the local trajectory, gradually turning the descent direction toward coordinate $i+1$ so that the selected step raises it farther above its activation level.

We shared this intuition with ChatGPT-6 Astra.
Through several rounds of substantive interaction and detailed calculations by ChatGPT-6 Astra Ultra, we designed a smooth convex hard function whose local gradient rotates along a circular arc and whose coefficient of $s_i$ approaches $1$.
Combining this construction with a refinement of the recursive analysis in \citet[Section~3]{JungChoYun2026} gave the present result.

\bibliographystyle{plainnat}
\bibliography{ref}

\clearpage
\appendix
\crefalias{section}{appendix}
\section{Proofs for the Hard-Function Construction}
\label{app:local}


\begin{fact}[Properties of Moreau envelope]
\label{lem:moreau-envelope}
Let $g:\R^d\to(-\infty,+\infty]$ be proper, closed, and convex.  Its Moreau
envelope with parameter $1$ and its proximal map are
\[
  \operatorname{env}_1 g(x)
  :=\min_{y\in\R^d}\left\{g(y)+\frac12\|x-y\|^2\right\},
  \qquad
  \operatorname{prox}_g(x)
  :=\argminop_{y\in\R^d}\left\{g(y)+\frac12\|x-y\|^2\right\}.
\]
The function $\operatorname{env}_1 g$ is finite, convex, and continuously differentiable, and its gradient is $1$-Lipschitz and satisfies $\nabla\operatorname{env}_1 g(x)=x-\operatorname{prox}_g(x)$~\citep[Propositions~5.b, 7.b, and~7.d]{Moreau1965}.

Moreover, if $K\subset\R^d$ is nonempty, closed, and convex and
$\sigma_K(x):=\sup_{v\in K}\langle v,x\rangle$ is its support function, then
\begin{equation*}
  \nabla\operatorname{env}_1\sigma_K(x)=\Pi_K(x),
\end{equation*}
where $\Pi_K$ denotes Euclidean projection onto $K$.
Finally, if $g\ge0$, then $\operatorname{env}_1g\ge0$ and
$\operatorname{env}_1g$ and $g$ have the same zero set. 
For more background on Moreau envelopes and proximal maps, see \citet[Sections~1.G and~2.D]{RockafellarWets1998} and \citet[Sections~3.1 and~6.4.2]{ParikhBoyd2014}.
\end{fact}


Fix $\epsilon,\gamma>0$ and recall that $c=\sqrt{1+2\gamma+\epsilon^2}$ and $R=\sqrt{\epsilon^2+(1+\gamma)^2}$.


\subsection{Proof of Lemma~\ref{lem:local-backward}}
\label{local:app-backward}
\begin{proof}
\label{local:backward-proof}

Set
\begin{equation}
  v(q):=(p(q),-q),\qquad
  p(q):=\sqrt{R^2-(q+\gamma)^2},\qquad
  0\le q\le1.
  \label{local:arc}
\end{equation}
The vectors $v(q)$ lie on the circle of radius $R$ centered at $(0,\gamma)$, and
\begin{equation}
  v(0)=(c,0),\qquad v(1)=(\epsilon,-1),\qquad
  \epsilon\le p(q)\le c,\qquad \frac{q}{p(q)}\le\frac1\epsilon.
  \label{local:arc-bounds}
\end{equation}
Let
\[
  K:=\operatorname{conv}\bigl(\{0\}\cup\{v(q):0\le q\le1\}\bigr).
\]
Consider stepsizes $\alpha_1,\ldots,\alpha_m\ge0$ and write their total stepsize as $S=\sum_{j=1}^m\alpha_j$.
Write $r_j=(u_j,-y_j)$ and $v_j=(p_j,-q_j)$ for $0\le j\le m$.
Start with $r_m=v_m=(\epsilon,-1)$.
For $j=m,m-1,\ldots,1$, define
\[
  \Delta_j:=y_j+\gamma-\alpha_j\gamma.
\]
If $\Delta_j\le u_j\gamma/c$, choose $v_{j-1}=(p_{j-1},-q_{j-1}):=(c,0)$.
Otherwise, choose
\begin{equation}
  D_j:=\sqrt{u_j^2+\Delta_j^2},\qquad
  q_{j-1}:=\frac{R\Delta_j}{D_j}-\gamma,\qquad
  p_{j-1}:=p(q_{j-1})=\frac{Ru_j}{D_j}.
  \label{local:backward-direction}
\end{equation}
In both cases define
\[
  r_{j-1}:=r_j+\alpha_jv_{j-1}
  =\bigl(u_j+\alpha_jp_{j-1},-y_j-\alpha_jq_{j-1}\bigr),
\]
so that
\[
  u_{j-1}=u_j+\alpha_jp_{j-1},\qquad
  y_{j-1}=y_j+\alpha_jq_{j-1}.
\]
The geometric relationship is illustrated in \cref{fig:backward-similar-triangles}.
Under this choice, we claim that for $0\le j\le m$,
\begin{equation}
  v_j=\Pi_K(r_j),\qquad
  u_j\ge\epsilon,\qquad y_j\ge1,\qquad \frac{y_j}{u_j}\le\frac1\epsilon.
  \label{local:projection-invariants}
\end{equation}
Furthermore, there is an index $j_\star\in\{0,\ldots,m\}$ such that
\begin{equation}\label{eq:j-star}
    \Delta_j>\frac{u_j\gamma}{c}
  \quad\Longleftrightarrow\quad
  j_\star<j\le m,
\end{equation}
where we take $j_\star=m$ if the strict inequality never holds.
\begin{figure}[H]
  \centering
  \includegraphics[width=\linewidth]{figures/backward-similar-triangles.pdf}
  \caption{An illustration of \eqref{local:backward-direction}. The shaded triangles are similar, with scale factor $\alpha_j$.}
  \label{fig:backward-similar-triangles}
\end{figure}

\paragraph{Proof of \eqref{local:projection-invariants} and \eqref{eq:j-star}.}
We prove \eqref{local:projection-invariants} by backward induction on \(j\). At \(j=m\), \(r_m=v_m\in K\), so \(v_m=\Pi_K(r_m)\), and the coordinate inequalities hold with equality. 
Fix \(j\in\{1,\ldots,m\}\) and assume the coordinate ineqaulities hold at \(r_j\). We verify \eqref{local:projection-invariants} at \(r_{j-1}\).

We first prove the coordinate inequalities. Suppose first that $\Delta_j>u_j\gamma/c$, so that \eqref{local:backward-direction} defines $v_{j-1}$.
By $\alpha_j\ge0$ and the inductive bounds $u_j\ge\epsilon$ and $y_j/u_j\le1/\epsilon$,
\[
  \frac\gamma c<\frac{\Delta_j}{u_j}
  \le\frac{y_j+\gamma}{u_j}
  =\frac{y_j}{u_j}+\frac\gamma{u_j}
  \le\frac1\epsilon+\frac\gamma\epsilon
  =\frac{1+\gamma}{\epsilon}.
\]
Define $\psi(s):=Rs/\sqrt{1+s^2}-\gamma$ for $s\ge0$.
Since $R^2=c^2+\gamma^2$ and $R^2=\epsilon^2+(1+\gamma)^2$,
\[
  \psi\!\left(\frac\gamma c\right)
  =\frac{R\gamma}{\sqrt{c^2+\gamma^2}}-\gamma=0,
  \qquad
  \psi\!\left(\frac{1+\gamma}{\epsilon}\right)
  =\frac{R(1+\gamma)}{\sqrt{\epsilon^2+(1+\gamma)^2}}-\gamma=1.
\]
Moreover, $\psi'(s)=R(1+s^2)^{-3/2}>0$.
Because $q_{j-1}=\psi(\Delta_j/u_j)$, the preceding bounds therefore give
\[
  0<q_{j-1}\le1.
\]
If instead $\Delta_j\le u_j\gamma/c$, then $(p_{j-1},q_{j-1})=(c,0)$.
Thus, in both branches, \eqref{local:arc-bounds} gives $p_{j-1}\ge\epsilon$ and $q_{j-1}\ge0$, so $u_{j-1}\ge\epsilon$ and $y_{j-1}\ge1$.
Moreover,
\[
  \frac{y_{j-1}}{u_{j-1}}
  =\frac{u_j}{u_j+\alpha_jp_{j-1}}\frac{y_j}{u_j}
   +\frac{\alpha_jp_{j-1}}{u_j+\alpha_jp_{j-1}}
    \frac{q_{j-1}}{p_{j-1}}
  \le\frac1\epsilon,
\]
as the two coefficients on the right sum to one, and \eqref{local:projection-invariants} and \eqref{local:arc-bounds} bound the two ratios by $1/\epsilon$.
This proves the coordinate inequalities for the induction step.

Now we prove that $v_{j-1}=\Pi_K(r_{j-1})$.
Suppose again that $\Delta_j>u_j\gamma/c$.
Then
\[
  v_{j-1}-(0,\gamma)
  =\frac{R}{D_j}(u_j,-\Delta_j),
  \qquad
  r_j=(u_j,-\Delta_j)+(1-\alpha_j)(0,\gamma).
\]
Using $r_{j-1}=r_j+\alpha_jv_{j-1}$, we obtain
\begin{equation}
\begin{aligned}
  r_{j-1}-v_{j-1}
  &=r_j+(\alpha_j-1)v_{j-1}\\
  &=\left(\frac{D_j}{R}+\alpha_j-1\right)
    \bigl(v_{j-1}-(0,\gamma)\bigr).
\end{aligned}
  \label{local:curved-normal}
\end{equation}
For every $\alpha_j\ge0$, the reverse triangle inequality and $u_j\ge\epsilon$, $y_j\ge1$ give
\[
\begin{aligned}
  D_j
  &=\bigl\|(u_j,y_j+\gamma)-\alpha_j(0,\gamma)\bigr\|\\
  &\ge\sqrt{u_j^2+(y_j+\gamma)^2}-\alpha_j\gamma\\
  &\ge\sqrt{\epsilon^2+(1+\gamma)^2}-\alpha_j\gamma
   =R-\alpha_j\gamma,
\end{aligned}
\]
hence
\[
  \frac{D_j}{R}+\alpha_j-1
  \ge\frac{R-\alpha_j\gamma}{R}+\alpha_j-1
  =\alpha_j\left(1-\frac\gamma R\right)\ge0.
\]
Therefore, the vector $v_{j-1}-(0,\gamma)$ is an outward normal to the circle, and \eqref{local:curved-normal} proves $v_{j-1}=\Pi_K(r_{j-1})$.
Now, suppose that $\Delta_j\le u_j\gamma/c$, so that $v_{j-1}=(c,0)$, and set $\mu_j:=u_j/c+\alpha_j-1$.
The defining inequality is equivalent to
\[
  \Delta_j=y_j+\gamma-\alpha_j\gamma\le\frac{u_j\gamma}{c}
  \quad\Longleftrightarrow\quad
  y_j\le\gamma\left(\frac{u_j}{c}+\alpha_j-1\right)=\gamma\mu_j.
\]
Since $y_j\ge1$, we have $\mu_j>0$, and
\begin{equation}
  r_{j-1}-(c,0)=(c\mu_j,-y_j)
  =\mu_j(c,-\gamma)+(0,\gamma\mu_j-y_j).
  \label{local:horizontal-normal}
\end{equation}
For every $w\in K$, the containment of $K$ in the disk centered at $(0,\gamma)$ and the inequality $w^{(2)}\le0$ give
\[
  \langle(c,-\gamma),w-(c,0)\rangle\le0,
  \qquad
  \langle(0,1),w-(c,0)\rangle=w^{(2)}\le0.
\]
Thus $(c,-\gamma)$ and $(0,1)$ belong to the normal cone of $K$ at $(c,0)$.
Since \eqref{local:horizontal-normal} is a nonnegative combination of these two vectors, it proves $v_{j-1}=\Pi_K(r_{j-1})$.

Finally, if $j\ge2$, then $y_{j-1}=y_j$ and $u_{j-1}=u_j+\alpha_jc$, so
\[
  \Delta_{j-1}
  =y_{j-1}+\gamma-\alpha_{j-1}\gamma
  \le y_j+\gamma
  \le\frac{u_j\gamma}{c}+\alpha_j\gamma
  =\frac{u_{j-1}\gamma}{c}.
\]
Hence, once $\Delta_j\le u_j\gamma/c$ holds at an index $j$, it also holds at every earlier index, and the construction chooses $v_i=(c,0)$ for all $0\le i<j$.
This completes the proof of \eqref{eq:j-star}.


\paragraph{Bounding the total stepsize for $j>j_\star$.}
\label{local:turning-proof}

Set $\sigma_j:=(y_j+\gamma)/u_j$ for $0\le j\le m$.
For every index $j$ satisfying $\Delta_j>u_j\gamma/c$, the recursion from $r_j$ to $r_{j-1}$ gives
\(\sigma_j-\sigma_{j-1}=\alpha_j\gamma/u_j,\) and thus
\begin{equation}
  \frac{u_{j-1}}{u_j}
  =1+\frac{\alpha_j}{u_j}p_{j-1}
  =1+\frac R\gamma
    \frac{\sigma_j-\sigma_{j-1}}
         {\sqrt{1+\sigma_{j-1}^2}}.
  \label{local:turning-recursion}
\end{equation}
For $A(\sigma)=\sigma+\sqrt{1+\sigma^2}$, convexity gives
\[
\begin{aligned}
  A(\sigma_j)
  &\ge A(\sigma_{j-1})
     +A'(\sigma_{j-1})(\sigma_j-\sigma_{j-1})\\
  &=A(\sigma_{j-1})\left(
     1+\frac{\sigma_j-\sigma_{j-1}}
              {\sqrt{1+\sigma_{j-1}^2}}
     \right),
\end{aligned}
\]
where we used $A'(\sigma)=A(\sigma)/\sqrt{1+\sigma^2}$.
Since $R/\gamma>1$, Bernoulli's inequality and \eqref{local:turning-recursion} give
\begin{equation}
  \frac{u_{j-1}}{u_j}
  =1+\frac R\gamma \frac{\sigma_j-\sigma_{j-1}}{\sqrt{1+\sigma_{j-1}^2}}
  \le\left(1+\frac{\sigma_j-\sigma_{j-1}}{\sqrt{1+\sigma_{j-1}^2}}\right)^{R/\gamma}
  \le\left(
    \frac{A(\sigma_j)}{A(\sigma_{j-1})}
  \right)^{R/\gamma}.
  \label{local:turning-telescope}
\end{equation}
Thus the ratios in \eqref{local:turning-telescope} telescope over $j=j_\star+1,\ldots,m$.
At the terminal point,
\[
  u_m=\epsilon,\qquad \sigma_m=\frac{1+\gamma}{\epsilon},
\]
and, for every $j_\star<j\le m$,
\[
  \sigma_{j-1}
  =\frac{y_j+\gamma-\alpha_j\gamma}{u_j}
  =\frac{\Delta_j}{u_j}
  >\frac\gamma c.
\]
Define
\[
  H_{\mathrm{turn}}
  :=\sum_{j:\,\Delta_j>u_j\gamma/c}\alpha_j
  =\sum_{j=j_\star+1}^m\alpha_j,
\]
then \eqref{local:arc-bounds} gives
\[
  u_{j_\star}
  =u_m+\sum_{j=j_\star+1}^m\alpha_jp_{j-1}
  \ge\epsilon(1+H_{\mathrm{turn}}).
\]
On the other hand, telescoping \eqref{local:turning-telescope} gives
\[
  \frac{u_{j_\star}}{\epsilon}
  \le\left[
    \frac{A(\sigma_m)}{A(\sigma_{j_\star})}
  \right]^{R/\gamma}.
\]
The endpoint values satisfy
\[
  A(\sigma_m)
  =A\!\left(\frac{1+\gamma}{\epsilon}\right)
  =\frac{1+\gamma+R}{\epsilon},
  \qquad
  A(\sigma_{j_\star})
  \ge A\!\left(\frac\gamma c\right)
  =\frac{\gamma+R}{c}.
\]
Consequently,
\[
  \frac{u_{j_\star}}{\epsilon}
  \le\left[
    \frac{c(1+\gamma+R)}{\epsilon(\gamma+R)}
  \right]^{R/\gamma}.
\]
It follows that
\begin{equation*}
  H_{\mathrm{turn}}\le H
  :=\left[\frac{c(1+\gamma+R)}{\epsilon(\gamma+R)}\right]^{R/\gamma}-1.
\end{equation*}
This conclusion also holds when $j_\star=m$, in which case $H_{\mathrm{turn}}=0$.

Finally, the displacements are
\begin{equation*}
  P=\sum_{j=1}^m\alpha_jp_{j-1},\qquad
  Q=\sum_{j=1}^m\alpha_jq_{j-1}.
\end{equation*}
When $\Delta_j\le u_j\gamma/c$, we have $q_{j-1}=0$; when $\Delta_j>u_j\gamma/c$, we have $0<q_{j-1}\le1$.
Consequently,
\begin{equation*}
  P\le cS,\qquad
  0\le Q
  =\sum_{j:\,\Delta_j>u_j\gamma/c}\alpha_jq_{j-1}
  \le H_{\mathrm{turn}}\le H.
\end{equation*}
\end{proof}

\subsection{Proof of Lemma~\ref{lem:local-transfer}}
\label{local:app-transfer}

\begin{proof}
Apply \Cref{lem:local-backward} to the prescribed stepsizes $\alpha_1,\ldots,\alpha_m$.
Retain its points $r_j$, vectors $v_j$, and the displacements $P,Q$ in \eqref{local:displacement-constants}, which satisfy \eqref{local:displacement-bounds}.
Define the constant
\begin{equation*}
    C_0:=\epsilon+\frac{H+1}{\epsilon},
\end{equation*}
then 
\begin{equation}\label{eq:cSplusC0}
    r_0=(\epsilon+P,-1-Q),\qquad
  \epsilon+P+\frac{1+Q}{\epsilon}\le cS+C_0.
\end{equation}

For the given final stepsize $b\ge1+\epsilon^{-2}$ and constant $A\ge C_0$, set
\begin{equation*}
  \rho:=\frac1{cS+\epsilon b+A},\qquad
  z_\circ:=(1,0)-\rho r_0.
\end{equation*}
Let $\mathcal V=\{(0,0),v_0,\ldots,v_m,(c,0)\}$ and define
\begin{equation}
  \phi(z):=\rho\max_{v\in\mathcal V}\langle v,z-z_\circ\rangle,
  \qquad f:=\operatorname{env}_1\phi.
  \label{local:finite-envelope}
\end{equation}
The finite convex hull $K_{\mathcal V}=\operatorname{conv}(\mathcal V)$ contains every $v_j$ and is contained in $K$.
Thus each $v_j$ remains the projection of $r_j$ onto the smaller set.
Since $(0,0)\in\mathcal V$, we have $\phi\ge0$; hence \Cref{lem:moreau-envelope} shows that $f$ is nonnegative, convex, and $1$-smooth.
Furthermore,
\[
  \phi(z)=\rho\sigma_{K_{\mathcal V}}(z-z_\circ)
  =\sigma_{\rho K_{\mathcal V}}(z-z_\circ),
  \qquad
  f(z)=\operatorname{env}_1\sigma_{\rho K_{\mathcal V}}(z-z_\circ),
\]
where the identity for $f$ follows by changing variables in the definition of
the envelope.  Therefore, \Cref{lem:moreau-envelope} and
$\Pi_{\rho K_{\mathcal V}}(\rho r)=\rho\Pi_{K_{\mathcal V}}(r)$ give
\begin{equation}
  \nabla f(z_\circ+\rho r_j)=\rho v_j\quad(0\le j\le m).
  \label{local:scaled-gradients}
\end{equation}
Let $U:=1-\rho(P+\epsilon b)$, $Y:=\rho Q$, and $G:=\rho b$, where $U\geq 0$ by \eqref{local:displacement-bounds}.
By \Cref{lem:local-backward}, if the updates on $f$ starts at $(1,0)$, it will end at $(U+\epsilon G,Y)$ just before the step of size $b$, where the gradient is $\rho(\epsilon,-1)$.
Therefore, the last step takes the iterate to $(U,Y+G)$.
During the intervening updates, the first coordinate decreases and the second increases as all stepsizes are nonnegative, and all gradients $(p,-q)$ satisfy $p>0$ and $q\ge0$.
This shows that the first coordinate remains at least $U$ through the final step, while the second lies in $[0,Y]$ before the final step.


For every $v=(p,-q)\in\mathcal V$, $q\le p/\epsilon$, $A\geq C_0$, and \eqref{eq:cSplusC0} gives
\[
  \langle v,-z_\circ\rangle
  =-p+\rho\bigl[p(\epsilon+P)+q(1+Q)\bigr]
  \le-p+\rho p(cS+C_0)\le0.
\]
Since $p\ge0$ and $q\ge0$, every affine piece in \eqref{local:finite-envelope} is nonpositive on $\{z_1\le0,z_2\ge0\}$.
On this region, $\phi=0$, so \Cref{lem:moreau-envelope} gives $f=0$.
Therefore
\begin{equation*}
  f(z_1,z_2)=0\qquad(z_1\le0,\ z_2\ge0).
\end{equation*}
Moreover, for every $v=(p,-q)\in\mathcal V$,
\[
  \langle v,(U,Y)-z_\circ\rangle
  =\rho\bigl[(1-b)\epsilon p+q\bigr]\le0,
\]
because $q\le p/\epsilon$ and $b\ge1+\epsilon^{-2}$.
Increasing the second coordinate decreases this inner product.
Thus $\phi=0$ on this region, and \Cref{lem:moreau-envelope} again gives
\begin{equation*}
  f(U,z_2)=0\qquad(z_2\ge Y).
\end{equation*}


\end{proof}

\subsection{Proof of Lemma~\ref{thm:transfer}}
\label{local:app-chain}

\begin{proof}
\label{local:chain-proof}
Recall from \Cref{sec:local} that
\begin{equation}
  B_0:=2(1+\epsilon^{-2}),\qquad
  a\ge 2C_0/c.
  \label{local:theorem-constants}
\end{equation}
Suppose first that $k\ge1$.
For each $i$, apply the local construction in \Cref{lem:local-transfer} with stepsizes $h_t/2$ for $t_{i-1}<t<t_i$, followed by the final step of size $b_i/2$, and with $A=ca/2$.
The assumed bound $b_i\ge B_0$, together with \eqref{local:theorem-constants}, verifies the hypotheses of \cref{lem:local-transfer}.
Write $f_i,U_i,Y_i$ for the resulting function and coordinates in \cref{lem:local-transfer}.
Use the distances $D_i$ and activation levels $\ell_i$ in \eqref{local:chain-parameters} and the two-coordinate components $\Phi_i$ in \eqref{local:chain-component}.
Each $\Phi_i$ is $1$-smooth on its pair of coordinates.
Use the one-sided Huber function $H_\delta$ from \cref{sec:overview}, with $\delta:=D_{k+1}/(1+s_{k+1})$, and define $F$ by \eqref{local:global-objective}.

\paragraph{Smoothness of $F$.}
The objective is convex and nonnegative.
All $\ell_i$ are nonnegative, so the first condition in \eqref{local:zero-regions} gives $\Phi_i(0)=0$, and the terminal term also vanishes at zero.
Thus $F(0)=\min F=0$.
To check smoothness, let $d=y-x$.
Summing the component smoothness inequalities gives
\begin{align*}
  F(y)&\le F(x)+\langle\nabla F(x),d\rangle
    +\frac14\left[
      \sum_{i=1}^k\bigl((d^{(i)})^2+(d^{(i+1)})^2\bigr)
      +(d^{(k+1)})^2\right]\\*
  &\le F(x)+\langle\nabla F(x),d\rangle+\frac12\|d\|^2.
\end{align*}
Each coordinate appears at most twice in the bracketed sum, proving that $F$ is $1$-smooth.

\paragraph{Analysis of trajectory.}
We now verify the trajectory starting from $x_0=e_1$.
At the start of block $i$, the claimed state is
\begin{equation}
  x_{t_{i-1}}^{(j)}=
  \begin{cases}
    \ell_j+D_jU_j,&j<i,\\
    \ell_i+D_i,&j=i,\\
    0,&j>i.
  \end{cases}
  \label{local:block-invariant}
\end{equation}
This holds initially.
Consider the path obtained from the trajectory of $f_i$ by the coordinate change in \eqref{local:chain-component}, leaving the other coordinates fixed.
At every iterate $t_{i-1}\le t<t_i$, it satisfies
\[
  x_t^{(i)}\ge\ell_i+D_iU_i\ge\ell_i,
  \qquad 0\le x_t^{(i+1)}\le D_iY_i=\ell_{i+1},
  \qquad x_t^{(j)}=0\quad(j\ge i+2).
\]
For every earlier two-coordinate component $j<i$, the first input to $f_j$ equals $U_j$ and the second is at least $\ell_{j+1}/D_j=Y_j$.
Hence the component remains zero by the second condition in \eqref{local:zero-regions}.
For every later two-coordinate component, the first input is nonpositive and the second is zero, so the component remains zero by the first condition in \eqref{local:zero-regions}.
The terminal Huber term is also zero before its activation.
A differentiable nonnegative function has zero gradient wherever it vanishes.
Thus the only nonzero gradient along this candidate path is that of $\Phi_i/2$.
By the scaling in \eqref{local:chain-component} and the factor $1/2$ in \eqref{local:global-objective}, each actual step $h_t$ induces exactly the local step $h_t/2$, so the candidate path is the GD trajectory.

The step of size $b_i$ gives
\[
  x_{t_i}^{(i)}=\ell_i+D_iU_i,\qquad
  x_{t_i}^{(i+1)}=D_i(Y_i+G_i)=\ell_{i+1}+D_{i+1},
\]
and all coordinates $j\ge i+2$ remain zero.
This proves \eqref{local:block-invariant} for the next block.
During that block, $x^{(i+1)}$ never drops below $\ell_{i+1}$, so the zero region of component $i$ continues to apply.
The induction verifies that the next component becomes active and every earlier component remains inactive.

\paragraph{Calculation of final loss.}
From $x_{t_k}$ onward, only the terminal Huber term is active.
The shifted last coordinate at $x_{t_k}$ equals $D_{k+1}$.
We prove by induction on $t$ that, for every $t_k\le t\le n$,
\[
\begin{aligned}
  x_t^{(k+1)}-\ell_{k+1}
  &=D_{k+1}-\frac{\delta}{2}\sum_{\tau=t_k+1}^{t}h_\tau\\
  &\ge D_{k+1}-\frac{\delta s_{k+1}}2
   =\delta\left(1+\frac{s_{k+1}}2\right)\ge\delta,
  \qquad \delta=\frac{D_{k+1}}{1+s_{k+1}}.
\end{aligned}
\]
At $t=t_k$, the sum is empty and the equality holds.
If it holds at $t-1$, then the terminal Huber term is in its affine region at
$x_{t-1}$, where its contribution to the last-coordinate gradient is
$\delta/2$; the update of size $h_t$ therefore gives the equality at $t$.
The inequalities use
$\sum_{\tau=t_k+1}^{t}h_\tau\le s_{k+1}$.
Thus this trajectory stays in the affine region and above its activation level.
Earlier components remain zero, and the final value is
\begin{equation}
  F(x_n)=\frac12\left[\delta\left(D_{k+1}-\frac{\delta s_{k+1}}{2}\right)
                             -\frac{\delta^2}{2}\right]
        =\frac{D_{k+1}^2}{4(1+s_{k+1})}.
  \label{local:terminal-value}
\end{equation}
Since $D_{k+1}=\prod_{i=1}^kG_i$ and $\|x_0-0\|=1$, the normalization in \eqref{eq:worst-case-error} gives $\mathcal R_n(h)\ge2F(x_n)$.
Dropping the factor $2$ proves \eqref{eq:transfer} for $k\ge1$.
For $k=0$, take $F(x)=H_\delta(x)/2$ in one dimension with
$x_0=1$, $\delta=1/(1+s_1)$, and minimizer zero.
The same calculation in \eqref{local:terminal-value} applies with $D_1=1$ and
proves the empty-product case as well.
\end{proof}

\section{Proofs for the Non-anytime Lower Bound}
\label{app:global}

We use the constants and conventions from \cref{sec:global}, with $0<\epsilon<1/2$, $c>\epsilon$, $a\ge\max\{2,2B_0+2,2C_0/c\}$, and $\lambda\ge a$.
For an empty subsequence, set $V_\mu(\varnothing)=\mu$.
A minimum over no indices is $+\infty$.
A checkpoint set selecting a zero step has $\Psi_\lambda(T;h)=+\infty$.
Increasing $a$ preserves the transfer bound.

\subsection[Proof of (\ref{global:eq-cost-lower-bound})]{Proof of \eqref{global:eq-cost-lower-bound}}
\label{global:app-cost-lower-bound}

\begin{lemma}
\label{lem:cost-lower-bound}
For every $n\ge0$, $h\in[0,\infty)^n$, and $\lambda\ge a$,
\begin{equation*}
  \mathcal R_n(h)\ge\frac{\lambda-1}{V_\lambda(h)^2}.
\end{equation*}
\end{lemma}

\begin{proof}[Proof of \cref{lem:cost-lower-bound}]
Since $\Psi_\lambda(\varnothing;h)<\infty$, a minimizing set cannot select a zero step.
Consider selecting a step of positive size $b$ at an index not already in the set.
Let $\ell$ be the sum of unselected steps preceding this step, starting after the preceding checkpoint if there is one.
Let $r$ be the sum of unselected steps following this step, ending before the next checkpoint if there is one.
If there is one, put $y=a+r+\epsilon b'/c$ when the next selected step has size $b'$; otherwise, put $y=\lambda+r$.
In both cases $y\ge a$, and insertion of this selected step of size $b$ multiplies $\Psi_\lambda$ by
\begin{equation}
  I(\ell,y;b)=\frac{[c(a+\ell)+\epsilon b]y}{b(\ell+b+y)}.
  \label{global:eq-insertion-ratio}
\end{equation}
The logarithmic derivatives are
\begin{align*}
  \partial_\ell\log I
    &=\frac{(c-\epsilon)b+c(y-a)}{[c(a+\ell)+\epsilon b](\ell+b+y)}>0,\\
  \partial_y\log I
    &=\frac{\ell+b}{y(\ell+b+y)}>0.
\end{align*}
If $0<b\le a/2$, then
\[
  I(\ell,y;b)\ge I(0,a;b)
  =\frac{(ca+\epsilon b)a}{b(a+b)}\ge\frac43.
\]
Deleting such a checkpoint strictly decreases $\Psi_\lambda$.
Thus every selected step of a minimizing set satisfies $b_i>a/2\ge B_0$, and \cref{thm:transfer} applies, including its empty-set case.
For a minimizing set $T$, it gives
\[
  \mathcal R_n(h)\ge
  \frac{(\lambda+s_{k+1})^2}
       {4(1+s_{k+1})V_\lambda(h)^2}.
\]
Finally, $ (\lambda+s)^2-4(\lambda-1)(1+s)=(s-\lambda+2)^2\ge0$ for every $s\ge0$, which proves the claim.
\end{proof}

\begin{remark}
The parameter $\lambda$ comes from the identity
\[
  \frac{1}{4(1+s)}
  =\max_{\lambda\ge2}\frac{\lambda-1}{(\lambda+s)^2},
  \qquad s\ge0.
\]
The expression $\lambda+s$ allows the factorization in \cref{global:lem-value-decomposition}.
We use $\lambda\ge a$ so that the same parameter range applies to each smaller problem.
\end{remark}

\subsection{Proof of Lemma~\ref{global:lem-value-decomposition}}
\label{global:app-value-decomposition}
\begin{proof}[Proof of \cref{global:lem-value-decomposition}]
Fix $t$ with $b=h_t>0$ and write $T=T_L\cup\{t\}\cup T_R$, where $T_L$ and $T_R$ contain the checkpoints before and after $t$.
Reindexing within each subsequence, the definition gives
\[
  \Psi_\lambda(T;h)
  =\frac cb\,
    \Psi_{a+\epsilon b/c}(T_L;h_{1:t-1})
    \Psi_\lambda(T_R;h_{t+1:n}).
\]
Indeed, the factor for $t$ is $\epsilon+c(a+s)/b=(c/b)(a+\epsilon b/c+s)$, where $s$ sums the stepsizes before $t$ and after the preceding checkpoint, or the start of the schedule.
Minimizing independently over $T_L$ and $T_R$ gives the term indexed by $t$ in \eqref{global:eq-value-decomposition}.
Every nonempty set with finite $\Psi_\lambda$ contains a positive step, so taking the minimum over $t$ and including the empty set proves the identity.
\end{proof}

\subsection{Proof of Lemma~\ref{thm:weighted-recursion}}
\label{global:app-uniform-growth}
\label{global:app-extremizers}
\label{global:app-gap-splitting}
\label{global:app-cost-preservation}
\label{global:app-weighted-splitting}

We adapt the log-submodularity and extremal-schedule arguments of \citet[Lemmas~C.1--C.2 and Appendix~C.6]{JungChoYun2026}.

For fixed $h$ and $\lambda$, call $T$ \emph{tight} if $\Psi_\lambda(T;h)=V_\lambda(h)$.
We use the same term for minimizers of $\widetilde\Psi_\Lambda$, defined in \eqref{global:eq-pure-cost}--\eqref{global:eq-pure-empty}.
We first show that $U_n(\lambda)$ is attained by a schedule for which the empty set and a nonempty checkpoint set are both tight.

\begin{lemma}[Log-submodularity]
\label{global:lem-log-submodularity}
For every positive schedule $h\in(0,\infty)^n$ and every two checkpoint sets $A,B$,
\begin{equation}
  \Psi_\lambda(A;h)\Psi_\lambda(B;h)
  \ge\Psi_\lambda(A\cap B;h)\Psi_\lambda(A\cup B;h).
  \label{global:eq-log-submodularity}
\end{equation}
The inequality is strict when $A$ and $B$ are disjoint and nonempty.
Consequently, tight sets are closed under unions and intersections, and two nonempty tight sets cannot be disjoint.
\end{lemma}

\begin{proof}[Proof of \cref{global:lem-log-submodularity}]
Consider inserting the same checkpoint into two sets, where the second contains the first.
Additional selected checkpoints decrease the sum $\ell$ in \eqref{global:eq-insertion-ratio}.
They also decrease the right-hand parameter $y$.
Indeed, selecting a closer step on the right, of size $b'$, replaces its full contribution to the sum of intervening stepsizes by $\epsilon b'/c$, and $\epsilon/c<1$.
If a checkpoint is added to the right where there was none, replacing $\lambda$ by $a$ can only decrease this parameter.
Thus both arguments of $I$ are no larger for the second set.

Add the elements of $B\setminus A$, one at a time, to each of $A\cap B$ and $A$.
At every insertion the multiplier for the second set is no larger.
Multiplying these inequalities gives \eqref{global:eq-log-submodularity}.
If $A$ and $B$ are disjoint and nonempty, the first insertion already compares the empty set with $A$.
Because $A$ contains a checkpoint, positivity of the schedule makes at least one of $\ell$ and $y$ strictly smaller for insertion into $A$ than for insertion into $\varnothing$.
This proves strictness.

If both $A$ and $B$ attain the minimum value $V:=V_\lambda(h)$, then
\[
  V^2\ge\Psi_\lambda(A\cap B;h)\Psi_\lambda(A\cup B;h)\ge V^2.
\]
Both sets on the right must therefore also be tight.
Strictness rules out disjoint nonempty tight sets.
\end{proof}

\begin{lemma}
\label{global:lem-extremal-schedule}
For every $n\ge1$ and $\lambda\ge a$, $U_n(\lambda)$ is finite, $U_n(\lambda)>U_{n-1}(\lambda)$, and the supremum defining $U_n(\lambda)$ is attained by a positive finite schedule.
At every maximizing schedule, the empty set and at least one nonempty checkpoint set are tight.
\end{lemma}

\begin{proof}[Proof of \cref{global:lem-extremal-schedule}]
Appending zero steps does not change $V_\lambda$, since those steps cannot be selected and contribute nothing to any sum of intervening stepsizes.
Hence $U_m(\lambda)\le U_j(\lambda)$ whenever $m\le j$.
We induct on $n$, starting with $U_0(\lambda)=\lambda$.

First, $U_n(\lambda)$ is finite.
For an arbitrary schedule, write $M=\max_t h_t$.
If $M\le a$, then $\Psi_\lambda(\varnothing;h)\le\lambda+na$.
If $M>a$, choose $j$ with $h_j=M$, select $j$ and no earlier checkpoint, and minimize over checkpoints in $h_{j+1:n}$.
This gives
\[
  V_\lambda(h)
  \le\left(\epsilon+\frac{c(a+\sum_{t<j}h_t)}{M}\right)
       U_{n-j}(\lambda)
  \le(\epsilon+cn)U_{n-1}(\lambda)<\infty.
\]
Here $\sum_{t<j}h_t\le(n-1)M$ and $a/M<1$.

Next, take a maximizing $(n-1)$-step schedule supplied by induction and append a positive step $t$.
For every checkpoint set that omits the new step, $\Psi_\lambda$ strictly increases, because its final sum of unselected stepsizes increases by $t$.
For every checkpoint set that selects the new step, $\Psi_\lambda$ tends to infinity as $t\downarrow0$.
There are finitely many checkpoint sets, so for small enough $t>0$ every such value exceeds $U_{n-1}(\lambda)$.
Consequently $U_n(\lambda)>U_{n-1}(\lambda)$.

To prove attainment, choose a maximizing sequence.
If it has an unbounded subsequence, pass to a further subsequence converging coordinatewise in the compact extended interval $[0,\infty]$ and having at least one infinite coordinate limit.
Let $j$ be the first such coordinate.
All coordinates before $j$ are bounded.
Selecting $j$ and minimizing over checkpoints in the suffix, as above, yields
\[
  \limsup V_\lambda(h)
  \le\epsilon U_{n-j}(\lambda)
  \le\epsilon U_{n-1}(\lambda)<U_n(\lambda),
\]
a contradiction.
Thus the maximizing sequence is bounded, and a subsequence converges to $\bar h\in[0,\infty)^n$.

Suppose $\bar h$ has a zero coordinate.
Delete all of its zero coordinates to obtain a schedule $\bar h^+$ of length $m<n$, and fix a minimizing checkpoint set for $\bar h^+$.
Use the corresponding indices in the converging sequence, omitting every coordinate with zero limit.
The selected steps have positive limits, and the contributions of omitted vanishing steps to sums of intervening stepsizes tend to zero.
The corresponding values of $\Psi_\lambda$ converge to $V_\lambda(\bar h^+)$.
Therefore
\[
  U_n(\lambda)\le V_\lambda(\bar h^+)
  \le U_m(\lambda)\le U_{n-1}(\lambda),
\]
again a contradiction.
The limit is positive.
On the positive orthant, $V_\lambda$ is the minimum of finitely many continuous functions, so it is continuous and $\bar h$ attains $U_n(\lambda)$.
The same argument after deleting zero steps shows that every maximizing schedule is positive.

It remains to establish tightness.
If the empty set were not tight, all tight sets would be nonempty.
Their intersection is nonempty and tight by \cref{global:lem-log-submodularity}.
Decreasing the stepsize at one index in their common intersection strictly increases $\Psi_\lambda$ at every initially tight set. That step has factor $\epsilon+c(a+s)/b$ and contributes to no sum of intervening stepsizes in those sets.
By continuity, the value at every other set remains strictly above the original minimum under a sufficiently small perturbation.
The minimum value would then increase, contradicting maximality.
Hence the empty set is tight.
If it were the only tight set, increasing any step slightly would increase $\Psi_\lambda(\varnothing;h)$ while preserving the slack of all other sets, giving the same contradiction.
A nonempty tight set therefore exists.
\end{proof}

We next prove the splitting inequality for the intervening updates, which is the key step in the proof of \cref{thm:weighted-recursion}.
We adapt the splitting argument in the proof of \citet[Lemma~3.2, Appendix~C.6]{JungChoYun2026}, retaining a correction term to handle the shifted parameters.

\begin{lemma}
\label{global:lem-gap-splitting}
Let $n\ge1$, $\lambda\ge a$, and let $h$ attain $U_n(\lambda)$.
Choose a nonempty tight set $T=\{t_1<\cdots<t_k\}$ and put $W=U_n(\lambda)$.
For $1\le i\le k+1$, let $h^{(i)}$ be the sequence of stepsizes indexed by $t_{i-1}<t<t_i$.
Then
\begin{equation}
  V_{a+\epsilon b_i/c}(h^{(i)})=a+s_i+\epsilon b_i/c\quad(1\le i\le k),
  \qquad
  V_\lambda(h^{(k+1)})=\lambda+s_{k+1}.
  \label{global:eq-gap-identities}
\end{equation}
If $\nu\in[1/2,1)$ satisfies \eqref{global:eq-scalar-hypothesis} with \eqref{global:eq-kappa-choice}, then
\begin{equation*}
  W^\nu+\sum_{i=1}^k(\epsilon b_i/c)^\nu
  \le(\lambda+s_{k+1})^\nu+\sum_{i=1}^k(a+s_i+\epsilon b_i/c)^\nu.
\end{equation*}
\end{lemma}

\begin{proof}[Proof of \cref{global:lem-gap-splitting}]
We first prove \eqref{global:eq-gap-identities} using tightness of $T$.
For any checkpoint set $S$ among the updates in $h^{(i)}$, with its indices understood in the original schedule, direct factorization gives
\[
  \frac{\Psi_\lambda(T\cup S;h)}{\Psi_\lambda(T;h)}
  =\frac{\Psi_{a+\epsilon b_i/c}(S;h^{(i)})}{a+s_i+\epsilon b_i/c}
  \qquad(1\le i\le k).
\]
The left side is at least one and the empty set of additional checkpoints attains equality.
The same argument for the updates after the last checkpoint replaces the right side by $\Psi_\lambda(S;h^{(k+1)})/(\lambda+s_{k+1})$.
This proves \eqref{global:eq-gap-identities}.

For the splitting inequality, we make a substitution and use two auxiliary lemmas, proved below.
For the selected stepsizes $b_i$ and total intervening stepsizes $s_i$, put
\begin{equation}
  B_i:=\frac{(c-\epsilon)b_i}{c},\qquad g_i:=s_i+\epsilon b_i/c,\qquad
  \Lambda:=\lambda+s_{k+1}.
  \label{global:eq-cost-move}
\end{equation}
These variables satisfy
\[
  g_i+B_i=s_i+b_i,
  \qquad
  \frac{(c-\epsilon)(a+g_i)}{B_i}=\frac{c(a+s_i)}{b_i}+\epsilon.
\]
When checkpoint $i$ is skipped, the first identity preserves its contribution to the sums of stepsizes.
When it is selected, the second preserves its factor in $\Psi_\lambda$.
This substitution changes only the expression for $\Psi_\lambda$, not the GD schedule.

For the splitting argument, allow arbitrary $g_1,\ldots,g_k\ge0$, $B_1,\ldots,B_k\ge0$, and $\Lambda\ge a$.
For a nonempty $S=\{i_1<\cdots<i_m\}\subseteq\{1,\ldots,k\}$, with $i_0=0$, set
\[
\begin{aligned}
  L_j(S)&:=\sum_{r=i_{j-1}+1}^{i_j}g_r
          +\sum_{r=i_{j-1}+1}^{i_j-1}B_r,\\
  Y(S)&:=\Lambda+\sum_{r>i_m}(g_r+B_r).
\end{aligned}
\]
Define the auxiliary function
\begin{align}
  \widetilde\Psi_\Lambda(S;g,B)
    &:=Y(S)\prod_{j=1}^m\frac{(c-\epsilon)(a+L_j(S))}{B_{i_j}},
    \label{global:eq-pure-cost}\\
  \widetilde\Psi_\Lambda(\varnothing;g,B)
    &:=\Lambda+\sum_{i=1}^k(g_i+B_i).
    \label{global:eq-pure-empty}
\end{align}
Set $\widetilde\Psi_\Lambda(S;g,B)=+\infty$ if $S$ contains an index with $B_i=0$.

By \cref{global:prop-cost-preservation}, the substitution preserves $\Psi_\lambda$ for every subset of $T$.
Since the empty set minimizes $\Psi_\lambda$ by \cref{global:lem-extremal-schedule}, it also minimizes $\widetilde\Psi_\Lambda$.
Applying \cref{lem:weighted-splitting} gives
\begin{equation*}
  W^\nu+\kappa\sum_{i=1}^k\bigl((c-\epsilon)b_i/c\bigr)^\nu
  \le(\lambda+s_{k+1})^\nu+
      \sum_{i=1}^k(a+s_i+\epsilon b_i/c)^\nu.
\end{equation*}
By \eqref{global:eq-kappa-choice}, $\kappa\bigl((c-\epsilon)b_i/c\bigr)^\nu=(\epsilon b_i/c)^\nu$.
This proves \eqref{global:eq-gap-splitting}.
\end{proof}

We now prove the two auxiliary lemmas used above.
For the first lemma, retain the schedule and variables in \eqref{global:eq-cost-move}.

\begin{lemma}
\label{global:prop-cost-preservation}
For every $S\subseteq\{1,\ldots,k\}$,
\begin{equation*}
  \widetilde\Psi_\Lambda(S;g,B)
  =\Psi_\lambda(\{t_i:i\in S\};h).
\end{equation*}
\end{lemma}


\begin{proof}[Proof of \cref{global:prop-cost-preservation}]
Skipping checkpoint $i$ leaves its contribution to a later sum of intervening stepsizes unchanged, since $g_i+B_i=s_i+b_i$.
If checkpoint $i_j$ is selected, let $\widetilde s_j:=\sum_{t_{i_{j-1}}<t<t_{i_j}}h_t$, with $t_{i_0}=0$.
Then $L_j(S)=\widetilde s_j+\epsilon b_{i_j}/c$.
The corresponding auxiliary factor is therefore
\[
  \frac{(c-\epsilon)(a+L_j(S))}{B_{i_j}}
  =\frac{c(a+\widetilde s_j+\epsilon b_{i_j}/c)}{b_{i_j}}
  =\frac{c(a+\widetilde s_j)}{b_{i_j}}+\epsilon.
\]
The final factors also agree.
For $S=\varnothing$, the equality follows directly from $g_i+B_i=s_i+b_i$.
\end{proof}

\begin{lemma}[Splitting with a correction term]
\label{lem:weighted-splitting}
Fix $\nu\in[1/2,1)$ and $\kappa\ge0$.
Suppose that, for every $x,y\ge a$ and $B>0$, the relations $w=x+y+B-a$ and $wB=(c-\epsilon) xy$ imply \eqref{global:eq-scalar-hypothesis}.
For any auxiliary problem \eqref{global:eq-pure-cost}--\eqref{global:eq-pure-empty} in which the empty set is a minimizer, write $W=\widetilde\Psi_\Lambda(\varnothing;g,B)$.
Then
\begin{equation}
  W^\nu+\kappa\sum_{i=1}^kB_i^\nu
  \le\Lambda^\nu+\sum_{i=1}^k(a+g_i)^\nu.
  \label{global:eq-weighted-splitting}
\end{equation}
\end{lemma}

With the choice of $\kappa$ in \eqref{global:eq-kappa-choice} and the substitution in \eqref{global:eq-cost-move}, the correction term $\kappa\sum_{i=1}^k B_i^\nu$ becomes $\sum_{i=1}^k(\epsilon b_i/c)^\nu$.
This is the term that cancels the extra contributions from the shifted parameters in the proof of \cref{thm:weighted-recursion}.

Fix $g,\Lambda$ and define
\begin{equation}
  J(B):=W^\nu+\kappa\sum_{i=1}^kB_i^\nu,
  \qquad W=\Lambda+\sum_{i=1}^k(g_i+B_i),
  \label{global:eq-weighted-objective}
\end{equation}

\begin{proof}[Proof of \cref{lem:weighted-splitting}]
We induct on $k$.
For $k=0$, $W=\Lambda$ and the assertion is equality.
For $k\ge1$, hold $g$ and $\Lambda$ fixed and maximize $J(B)$ from \eqref{global:eq-weighted-objective} over $B_i\ge0$ for which the empty set minimizes $\widetilde\Psi_\Lambda$.
It suffices to prove the desired bound at a maximizer, since its right side is fixed.

\paragraph{Compactness and zero coordinates.}
The feasible set is nonempty because $B=0$ is feasible.
For the singleton constraint at $i$, define
\[
  x_i=a+\sum_{r\le i}g_r+\sum_{r<i}B_r,
  \qquad y_i=\Lambda+\sum_{r>i}(g_r+B_r).
\]
Its constraint is $WB_i\le(c-\epsilon) x_i y_i$.
Since $W\ge y_i>0$, we have $B_i\le(c-\epsilon) x_i$.
The latter bound depends only on earlier variables and fixed values of $g_i$, so applying it successively bounds every $B_i$.
For each nonempty $S$, multiplying $W\le\widetilde\Psi_\Lambda(S;g,B)$ by $\prod_{i\in S}B_i$, when all selected sizes are positive, gives a continuous polynomial inequality.
If a selected size is zero, the left side becomes zero and the right side is positive, agreeing with the $+\infty$ convention for $\widetilde\Psi_\Lambda$.
The feasible set is therefore closed and bounded, and the continuous function $J$ attains a maximum.

If a maximizing vector has $B_i=0$ with $i<k$, delete that checkpoint and replace $g_{i+1}$ by $g_i+g_{i+1}$.
This preserves the values of $\widetilde\Psi_\Lambda$ for all sets that omit $i$, and hence preserves empty-set optimality.
The induction hypothesis, followed by
\[
  (a+g_i+g_{i+1})^\nu
  \le(a+g_i)^\nu+(a+g_{i+1})^\nu,
\]
gives the desired bound.
If $B_k=0$, delete it and replace $\Lambda$ by $\Lambda+g_k$. Then use $(\Lambda+g_k)^\nu\le\Lambda^\nu+(a+g_k)^\nu$.
We may therefore assume that a maximizing vector has every $B_i>0$.

\paragraph{Finding a tight set with one checkpoint.}
At least one nonempty subset is tight.
Otherwise, a sufficiently small increase of any $B_i$ would preserve all constraints and strictly increase $J$.
The auxiliary expression $\widetilde\Psi_\Lambda$ has the same log-submodularity and strictness property as in \cref{global:lem-log-submodularity}. Its insertion ratio is
\[
  \frac{(c-\epsilon)(a+\ell)y}{B(\ell+B+y)},
\]
which is strictly increasing in $\ell$ and $y$ for $y\ge a$ and $B>0$.
The proof by successive insertions applies unchanged.
Choose an inclusion-minimal nonempty tight set $T_0$.
Every nonempty tight set intersects $T_0$, and that intersection is tight.
Minimality therefore implies that $T_0$ lies in every nonempty tight set.

Suppose first that $T_0$ contains two unequal sizes $u<v$, and put $r=v/u>1$.
Because both checkpoints are selected in every nonempty tight set, $\widetilde\Psi_\Lambda$ for each such set depends on them only through the reciprocal product $1/(uv)$.
Define
\[
  A_*:=\frac{v(W+u)}{u(W+v)},\qquad
  D_*:=\frac{W^{\nu-1}+\kappa u^{\nu-1}}
             {W^{\nu-1}+\kappa v^{\nu-1}}.
\]
The positive term $\Lambda$ gives $W>u+v$.
Since $\nu\ge1/2$ and $\kappa\ge0$,
\begin{equation*}
  A_*>\frac{r(r+2)}{2r+1}\ge\sqrt r
  \ge r^{1-\nu}\ge D_*\ge1.
\end{equation*}
For completeness, the first inequality follows because $r(z+1)/(z+r)$ increases in $z=W/u>1+r$.
The second follows, on writing $q=\sqrt r$, from $q^3-2q^2+2q-1=(q-1)(q^2-q+1)\ge0$.
The bound on $D_*$ follows by adding the same positive number $W^{\nu-1}$ to the numerator and denominator of a ratio at most $u^{\nu-1}/v^{\nu-1}=r^{1-\nu}$.

Choose $\theta\in(D_*,A_*)$ and perturb $u\mapsto u-t$, $v\mapsto v+\theta t$.
For every nonempty tight set, $\widetilde\Psi_\Lambda$ initially equals $W$.
The derivative of the difference between its value and the empty-set value is
\[
  W\left(\frac1u-\frac\theta v\right)-(\theta-1)>0,
\]
where the inequality is equivalent to $\theta<A_*$.
Meanwhile the objective has derivative
\[
  \nu\bigl[(\theta-1)W^{\nu-1}
     +\kappa(\theta v^{\nu-1}-u^{\nu-1})\bigr]>0,
\]
since $\theta>D_*$.
Positivity of the $B_i$ and the strictness of all other constraints persist for sufficiently small $t>0$.
This contradicts maximality.

If $T_0$ has at least two elements and all of its sizes are equal to $b$, perturb two of them to $b-t$ and $b+t+t^2/(2b)$.
Their product becomes $b^2-t^2/2-t^3/(2b)$, while $W$ increases by $t^2/(2b)$.
Thus every nonempty tight constraint gains slack
\[
  \frac{W-b}{2b^2}t^2+O(t^3)>0.
\]
The change in $J$ is
\[
  \frac{\nu t^2}{2b}
  \bigl[W^{\nu-1}+\kappa(2\nu-1)b^{\nu-1}\bigr]+O(t^3)>0.
\]
The leading coefficient is positive even for $\nu=1/2$ or $\kappa=0$, because $W^{\nu-1}>0$.
Again the remaining constraints retain their strict slack for small $t$.
This is a contradiction.
It follows that $T_0=\{r\}$ for some index $r$.

\paragraph{Splitting at that checkpoint.}
Write
\[
  x:=a+\sum_{i\le r}g_i+\sum_{i<r}B_i,
  \qquad y:=\Lambda+\sum_{i>r}(g_i+B_i).
\]
Tightness of $\{r\}$ gives $W=x+y+B_r-a$ and $WB_r=(c-\epsilon) xy$, with $x,y\ge a$.
Define the left and right auxiliary expressions by
\[
  \widetilde\Psi_L(S_L):=\widetilde\Psi_{a+g_r}(S_L;g_{1:r-1},B_{1:r-1}),
  \qquad
  \widetilde\Psi_R(S_R):=\widetilde\Psi_\Lambda(S_R;g_{r+1:k},B_{r+1:k}),
\]
where $S_L\subseteq\{1,\ldots,r-1\}$ and $S_R\subseteq\{r+1,\ldots,k\}$ use the original indices; in $\widetilde\Psi_R(S_R)$, each index is reduced by $r$.
For every such pair,
\begin{equation}
  \widetilde\Psi_\Lambda(S_L\cup\{r\}\cup S_R;g,B)
  =\frac{c-\epsilon}{B_r}\widetilde\Psi_L(S_L)\widetilde\Psi_R(S_R).
  \label{global:eq-pure-factorization}
\end{equation}
Their empty-set values are $\widetilde\Psi_L(\varnothing)=x$ and $\widetilde\Psi_R(\varnothing)=y$.
Taking $S_R=\varnothing$ in \eqref{global:eq-pure-factorization} and using optimality of the empty set in the full auxiliary problem gives $((c-\epsilon) y/B_r)\widetilde\Psi_L(S_L)\ge W=((c-\epsilon) y/B_r)x$.
Thus the empty set minimizes the left problem.
Taking $S_L=\varnothing$ gives the same conclusion for the right problem.

The induction hypothesis applies to both sides, including a side with no checkpoints, and yields
\begin{align*}
  x^\nu+\kappa\sum_{i<r}B_i^\nu
    &\le\sum_{i\le r}(a+g_i)^\nu,\\
  y^\nu+\kappa\sum_{i>r}B_i^\nu
    &\le\Lambda^\nu+\sum_{i>r}(a+g_i)^\nu.
\end{align*}
The scalar hypothesis gives $W^\nu+\kappa B_r^\nu\le x^\nu+y^\nu$.
Combining these three inequalities proves \eqref{global:eq-weighted-splitting}.
\end{proof}

With \cref{global:lem-gap-splitting} established, we return to the proof of \cref{thm:weighted-recursion}.

\begin{proof}[Proof of \cref{thm:weighted-recursion}]
We induct on $n$, simultaneously for all $\lambda\ge a$.
The case $n=0$ follows from $U_0(\lambda)=\lambda$.
For $n\ge1$, choose a maximizing schedule, put $W=U_n(\lambda)$, and choose a nonempty tight set with $k$ checkpoints.
By \cref{global:lem-gap-splitting}, we have \eqref{global:eq-gap-splitting}.
Let $m_i$ be the number of updates in $h^{(i)}$ for $1\le i\le k+1$.
Since $k\ge1$, each $m_i<n$, and $\sum_{i=1}^{k+1}m_i=n-k$.
The identities \eqref{global:eq-gap-identities} and the induction hypothesis at parameters $a+\epsilon b_i/c\ge a$ and $\lambda\ge a$ give
\begin{equation*}
\begin{aligned}
  (a+s_i+\epsilon b_i/c)^\nu&\le a^\nu m_i+(a+\epsilon b_i/c)^\nu
      &&(1\le i\le k),\\
  (\lambda+s_{k+1})^\nu&\le a^\nu m_{k+1}+\lambda^\nu.
\end{aligned}
\end{equation*}
Using $(a+\epsilon b_i/c)^\nu\le a^\nu+(\epsilon b_i/c)^\nu$ in \eqref{global:eq-gap-splitting} yields
\[
  W^\nu+\sum_{i=1}^k(\epsilon b_i/c)^\nu
  \le a^\nu n+\lambda^\nu+\sum_{i=1}^k(\epsilon b_i/c)^\nu.
\]
Cancelling the checkpoint terms proves \eqref{eq:cost-growth}.
\end{proof}

For $n\ge1$, let $p=1/\nu$ and take $\lambda=an^p$.
\cref{thm:weighted-recursion} gives $V_\lambda(h)\le U_n(\lambda)\le 2^p a n^p$.
Since $a\ge2$, we have $an^p-1\ge an^p/2$, and \cref{lem:cost-lower-bound} implies \eqref{eq:fixed-p-rate}.
It remains to verify the scalar hypothesis for $p$ arbitrarily close to $\psil$ and to quantify the dependence of $a$ on that choice.

\subsection{Proof of Lemma~\ref{global:lem-scalar-condition}}
\label{global:app-scalar}

For $x,y,B,w$ in the scalar hypothesis, set $u=x/w$, $v=y/w$, $p=1/\nu$, and $\tau=\epsilon^{1/p}$.
Here $0<u,v<1$, since $x,y\ge a$ and $B>0$.
The scalar relations imply
\begin{equation}
  u+v+(c-\epsilon) uv=1+\frac aw\ge1.
  \label{global:eq-normalized-scalar-constraint}
\end{equation}
Since $B/w=(c-\epsilon) uv$ and $\kappa=(\epsilon/(c-\epsilon))^{1/p}$, dividing \eqref{global:eq-scalar-hypothesis} by $w^\nu$ gives the target
\begin{equation}
  u^{1/p}+v^{1/p}-\tau(uv)^{1/p}\ge1.
  \label{global:eq-normalized-scalar-target}
\end{equation}
To see why the silver inequality appears, consider the ideal limiting scalar problem $c=1$, $\epsilon=0$.
Writing $X=u^{1/p}$ and $Y=v^{1/p}$, a violation would have $Y<1-X$.
The expression $X^p+Y^p+X^pY^p$ increases with $Y$, and its boundary value is
\[
  X^p+(1-X)^p+[X(1-X)]^p.
\]
If this expression is at most $1$, the target inequality \eqref{global:eq-normalized-scalar-target} must hold. Otherwise, \eqref{global:eq-normalized-scalar-constraint} would be violated.
By \cref{global:lem-silver-inequality}, $X^p+(1-X)^p+[X(1-X)]^p\le1$ at $p=\psil$.

\begin{lemma}[Silver inequality]
\label{global:lem-silver-inequality}
For every \(z\in[0,1]\),
\begin{equation*}
 z^{\psil}+(1-z)^{\psil}+[z(1-z)]^{\psil}\le 1.
\end{equation*}
Equality holds at \(z=0,1/2,1\).
\end{lemma}

\begin{proof}[Proof of \cref{global:lem-silver-inequality}]
Write $\alpha=\psil-1\in(0,1/3)$, and denote the left side of \eqref{global:eq-silver-inequality} by $F(z)$.
For $0<z<1/2$,
\[
  \frac{F'(z)}{\psil[z(1-z)]^\alpha}
  =\phi(z):=(1-z)^{-\alpha}-z^{-\alpha}+1-2z.
\]
Moreover,
\[
  \phi''(z)=\alpha(\alpha+1)
  \bigl[(1-z)^{-\alpha-2}-z^{-\alpha-2}\bigr]<0.
\]
The endpoint values satisfy
\[
  \phi(0+)=-\infty,\qquad \phi'(0+)=+\infty,\qquad
  \phi(1/2)=0,
\]
and
\[
  \phi'(1/2)=\alpha 2^{\alpha+2}-2
  <\frac13 2^{7/3}-2<0.
\]
Thus $\phi'$ has exactly one zero in $(0,1/2)$, and $\phi$ first increases and then decreases.
Since $\phi'(1/2)<0$, $\phi$ is positive immediately to the left of $1/2$.
It therefore has exactly one zero in $(0,1/2)$, with negative sign before that zero and positive sign after it.
Consequently, $F$ first decreases and then increases on $[0,1/2]$, and its maximum occurs at an endpoint.
With $\rho=1+\sqrt2$, those endpoint values are
\[
  F(0)=1,\qquad F(1/2)=\frac2\rho+\frac1{\rho^2}=1.
\]
Symmetry under $z\mapsto1-z$ proves the result.
\end{proof}

Fix $c,\epsilon$ with $0<c-\epsilon\le2$.
For $p\in[\psil,3/2]$, set
\[
 G_p(z):=z^p+(1-z)^p+(c-\epsilon)[z(1-z)]^p.
\]
Writing \(t_+=\max\{t,0\}\), we have
\begin{equation}
\label{global:eq-silver-slack}
 G_p(z) \le 1-\bigl[p-\psil-(c-\epsilon-1)_+\bigr]z(1-z), \qquad z\in[0,1].
\end{equation}

\begin{proof}[Proof of the slack estimate \eqref{global:eq-silver-slack}]
For $0<z<1$, the inequalities $-\log z\ge1-z$ and $-\log(1-z)\ge z$ give
\begin{align*}
  -\partial_pG_p(z)
  &\ge z^p(1-z)+(1-z)^p z\\
  &=z(1-z)\bigl[z^{p-1}+(1-z)^{p-1}\bigr]
  \ge z(1-z).
\end{align*}
The last inequality uses $0<p-1\le1/2$.
By \cref{global:lem-silver-inequality} and $[z(1-z)]^{\psil}\le z(1-z)$,
\[
  G_{\psil}(z)\le1+(c-\epsilon-1)_+z(1-z).
\]
Integrating the derivative bound from $\psil$ to $p$ proves the estimate.
Continuity gives the endpoint cases.
\end{proof}

\begin{proof}[Proof of \cref{global:lem-scalar-condition}]
Let $x,y\ge a$ and $B>0$ satisfy
\[
  w=x+y+B-a,\qquad wB=(c-\epsilon) xy.
\]
Set $u=x/w$ and $v=y/w$.
Since $a>0$, $x,y\ge a$, and $B>0$, we have $0<u,v<1$.
The two relations imply \eqref{global:eq-normalized-scalar-constraint}.
Dividing $w^\nu+\kappa B^\nu\le x^\nu+y^\nu$ by $w^\nu$, and using $B/w=(c-\epsilon) uv$, reduces the desired conclusion to \eqref{global:eq-normalized-scalar-target}.

Suppose this inequality fails.
Put $X=u^{1/p}$ and $Y=v^{1/p}$.
Then $0<X,Y<1$ and
\[
  Y<\frac{1-X}{1-\tau X}=:A.
\]
Here $0<A<1$.
Since the derivative of $t^p$ on $[0,1]$ is at most $p$, the mean value theorem gives
\begin{align*}
  \bigl[A^p-(1-X)^p\bigr](1+(c-\epsilon) X^p)
  &\le p\bigl[A-(1-X)\bigr](1+c-\epsilon)\\
  &\le\frac{p\tau(1+c-\epsilon)}{1-\tau}X(1-X)
  \le6\tau X(1-X).
\end{align*}
The final step uses $p\le3/2$, $(c-\epsilon)\le2$, and $\tau\le1/10$.
Consequently,
\begin{align*}
  u+v+(c-\epsilon) uv
  &<X^p+A^p+(c-\epsilon) X^pA^p\\
  &\le G_p(X)+6\tau X(1-X)\le1,
\end{align*}
where the last inequality follows from \eqref{global:eq-silver-slack} and \eqref{global:eq-scalar-sufficient}.
This contradicts \eqref{global:eq-normalized-scalar-constraint}.
\end{proof}

\subsection[Proof of (\ref{eq:parameter-growth})]{Proof of \eqref{eq:parameter-growth}}
\label{global:app-parameters}
\label{global:app-nonanytime}

Fix \(0<\eta\le1/10\) and use the parameter choice \eqref{global:eq-parameter-choice}.
These choices give \(p\in[\psil,3/2]\), \(0<\epsilon<1/2\), and \(\epsilon^{1/p}=\eta/16\le1/10\).
Also, \(\epsilon^2\le(\eta/16)^2\le\eta/6\), so
\[
 c^2=1+\frac{\eta}{2}+\epsilon^2 \le1+\frac{2\eta}{3} \le\left(1+\frac{\eta}{3}\right)^2.
\]
Consequently, \(0<(c-\epsilon)\le c\le1+\eta/3<2\), and
\[
 (c-\epsilon-1)_++6\epsilon^{1/p} \le\frac{\eta}{3}+\frac{3\eta}{8} =\frac{17\eta}{24}<\eta=p-\psil.
\]
\cref{global:lem-scalar-condition} therefore verifies \eqref{global:eq-scalar-hypothesis} with the coefficient \(\kappa=(\epsilon/(c-\epsilon))^{1/p}\) prescribed in \eqref{global:eq-kappa-choice}.
The bound in \cref{thm:weighted-recursion}, together with \eqref{eq:fixed-p-rate}, now gives the lower bound with exponent \(\psil+\eta\).
Taking \(\eta=1/10\) already gives \(\mathcal R_n(h)=\Omega(n^{-(p_{\mathrm{sil}}+1/10)})\), which also implies \(\mathcal R_n(h)=\Omega(n^{-q})\) for every \(q>p_{\mathrm{sil}}+1/10\).

To quantify the dependence on \(\eta\), take the constants from the transfer construction to be
\begin{equation*}
 \begin{gathered}
 R:=\sqrt{\epsilon^2+(1+\gamma)^2},\qquad H:=\left[\frac{c(1+\gamma+R)}{\epsilon(\gamma+R)}\right]^{R/\gamma}-1,\\
 C_0:=\epsilon+\epsilon^{-1}+H/\epsilon,\qquad B_0:=2(1+\epsilon^{-2}),\\
 a:=\max\{2,2C_0/c,2B_0+2\}.
 \end{gathered}
\end{equation*}
In particular, \(a\ge2B_0+2\), as required.
Under \eqref{global:eq-parameter-choice}, the quantities \(c\) and \(R\) are bounded above and bounded away from zero by absolute constants, while
\[
 \log(1/\epsilon)=p\log(16/\eta)=O(\log(1/\eta)), \qquad \frac{R}{\gamma}=O(1/\eta).
\]
The logarithm of the base defining \(1+H\) is \(O(\log(1/\eta))\).
Hence
\[
 \log(1+H) =\frac{R}{\gamma} \log\!\left(\frac{c(1+\gamma+R)}{\epsilon(\gamma+R)}\right) =O\!\left(\eta^{-1}\log(1/\eta)\right).
\]
Since \(C_0=\epsilon+(1+H)/\epsilon\) and \(\log B_0=O(\log(1/\eta))\), the choice of \(a\) gives \eqref{eq:parameter-growth} with an absolute implied constant. 

\section{Proofs for the Anytime Lower Bound}
\label{app:anytime}

\subsection{Proof of Lemma~\ref{any:lem:record-sum}}
\label{any:app-record-sum}

We prove \cref{any:lem:record-sum}.
Fix $p\in(\psil,\psil+1/10]$ and use the parameters $a,c,\epsilon,B_0$ from \cref{global:app-parameters}, with $\eta=p-\psil$.
Recall that $\nu=1/p$ and $0<\epsilon/c<1$.
We may take
\[
    r_{0,p}:=\min\left\{\frac18,\frac{1}{16c^2}\right\}.
\]
At a record time satisfying the lemma's assumptions, write
\[
    M:=M_n=h_n,
    \qquad R:=r_n,
    \qquad \xi:=H_{n-1},
    \qquad \Lambda:=a+\epsilon M/c,
    \qquad V:=V_\Lambda(\xi).
\]

\paragraph{A lower bound on $V$.}
The deletion argument in the proof of \cref{lem:cost-lower-bound} shows that deleting any selected step of size at most $a/2$ decreases $\Psi_\Lambda$.
Thus a set attaining $V$ selects only steps larger than $a/2\ge B_0$.
Append the final step $M$ to this set.
Let $s$ be the sum of the stepsizes after the last selected step in the prefix, or of all stepsizes in the prefix if no step is selected. The factor for the appended step becomes $c(\Lambda+s)/M$, since $c(a+s)+\epsilon M=c(\Lambda+s)$.
There are no updates after the selected final step, so \cref{thm:transfer} gives
\begin{equation}
    R\ge\frac{M^2}{4c^2V^2},
    \qquad
    V\ge\frac{M}{2c\sqrt R}\ge2M\ge2\epsilon M/c.
    \label{any:eq:prefix-cost-lower}
\end{equation}

\paragraph{Counting steps above a threshold.}
For $0<u\le M$, let
\[
    N_\xi(u):=\bigl|\{t<n:h_t>u\}\bigr|=k.
\]
List these $k$ steps in the order they occur as $b_1,\ldots,b_k$, and let $\xi^{(0)},\ldots,\xi^{(k)}$ be the blocks between them, including the initial and final blocks.
Write $m_j$ for the length of $\xi^{(j)}$.
Then $\sum_{j=0}^k(m_j+1)=n$.

Select $b_1,\ldots,b_k$ together with checkpoint sets attaining $V_{a+\epsilon b_i/c}(\xi^{(i-1)})$ for $1\le i\le k$ and $V_\Lambda(\xi^{(k)})$ for the final block.
Evaluating $\Psi_\Lambda$ at this set gives
\begin{equation}
    V_{a+\epsilon M/c}(\xi)
    \le V_{a+\epsilon M/c}(\xi^{(k)})
    \prod_{i=1}^k\frac{c}{b_i}
    V_{a+\epsilon b_i/c}(\xi^{(i-1)}).
    \label{any:eq:block-factorization}
\end{equation}
Indeed, multiplying the preceding block's final factor $a+s+\epsilon b_i/c$ by $c/b_i$ gives $\epsilon+c(a+s)/b_i$, the factor for $b_i$.
This factorization uses only the definition of $\Psi_\lambda$, so the selected $b_i$ need not exceed the threshold $B_0$ required by \cref{thm:transfer}.

Applying \eqref{eq:cost-growth} to each block and using $(a+\epsilon b/c)^\nu\le a^\nu+(\epsilon b/c)^\nu$ gives
\begin{align*}
    \frac{c}{b_i}V_{a+\epsilon b_i/c}(\xi^{(i-1)})
    &\le \epsilon\left(
        1+\frac{(ac/\epsilon)^\nu(m_{i-1}+1)}{b_i^\nu}
    \right)^{1/\nu},
    \\
    V_{a+\epsilon M/c}(\xi^{(k)})
    &\le\frac{\epsilon M}{c}\left(
        1+\frac{(ac/\epsilon)^\nu(m_k+1)}{M^\nu}
    \right)^{1/\nu}.
\end{align*}
Since $b_i>u$, $M\ge u$, and $1+z\le e^z$, substituting these estimates into \eqref{any:eq:block-factorization} yields
\[
    \frac{cV}{\epsilon M}
    \le\epsilon^k
    \exp\!\left\{
        \frac{(ac/\epsilon)^\nu}{\nu}\,n u^{-\nu}
    \right\}.
\]
Together with \eqref{any:eq:prefix-cost-lower}, this implies
\begin{equation*}
    N_\xi(u)
    \le\frac{(ac/\epsilon)^\nu}{\nu\log(1/\epsilon)}
    \,n u^{-\nu}.
\end{equation*}
Integrating this bound over the threshold gives
\begin{equation}
    \sum_{t=1}^{n-1}h_t
    =\int_0^M N_\xi(u)\,du
    \le\frac{(ac/\epsilon)^\nu}{\nu(1-\nu)\log(1/\epsilon)}
    \,nM^{1-\nu}.
    \label{any:eq:prefix-sum}
\end{equation}

\paragraph{Including the last step.}
By \eqref{any:eq:prefix-cost-lower} and \eqref{eq:cost-growth},
\[
    (2c\sqrt R)^{-\nu}M^\nu
    \le V^\nu
    \le a^\nu(n-1)+(a+\epsilon M/c)^\nu
    \le a^\nu n+(\epsilon M/c)^\nu.
\]
Our choice of $r_{0,p}$ ensures $(2c\sqrt R)^{-\nu}\ge2^\nu$, and hence
\[
    M^\nu\le\frac{a^\nu}{2^\nu-(\epsilon/c)^\nu}\,n.
\]
Adding $M\le a^\nu nM^{1-\nu}/(2^\nu-(\epsilon/c)^\nu)$ to \eqref{any:eq:prefix-sum} proves the lemma with
\begin{equation*}
    C_p=\frac{(ac/\epsilon)^\nu}{\nu(1-\nu)\log(1/\epsilon)}
        +\frac{a^\nu}{2^\nu-(\epsilon/c)^\nu}.
\end{equation*}
\hfill$\square$

The explicit choices of $r_{0,p}$ and $C_p$ also justify the dependence on $p$ used in \eqref{any:eq:parameter-growth}.
For $\delta=p-\psil\downarrow0$, both $\nu$ and $1-\nu$ stay bounded away from zero, while $c$ stays bounded and $\log(c/\epsilon)=O(\log(1/\delta))$.
Thus \eqref{eq:parameter-growth} gives $\log C_p=O(\delta^{-1}\log(1/\delta))$.
The additional constants in the derivation from \eqref{any:eq:final-step-lower} to \eqref{any:eq:fixed-p-rate} are bounded by fixed powers of $a,c,B_0$ and $C_p$, which gives the stated bound on $\log(1/c_p)$.

\subsection{Proof of Theorem~\ref{thm:anytime}}
\label{any:app-main-proof}
In this proof, $K_p$ denotes a positive constant depending only on $p$, whose value may change between occurrences.
Selecting no checkpoints in \cref{thm:transfer} gives
\begin{equation}
    r_n\ge\frac{1}{4(1+S_n)}.
    \label{any:eq:total-sum-lower}
\end{equation}
If $H$ is bounded, then $S_n=O(n)$, and this is already stronger than the desired bound.
We may therefore assume that $H$ is unbounded and consider its infinitely many strict record times $n\ge2$, where $h_n>M_{n-1}$.

At a record time with $M_n\ge B_0$, selecting only the final step in \cref{thm:transfer} gives
\begin{equation}
    r_n\ge\frac14 \left[\frac{M_n}{c(a+S_{n-1})+\epsilon M_n}\right]^2.
    \label{any:eq:final-step-lower}
\end{equation}
For $r_n\le r_{0,p}$, after decreasing $r_{0,p}$ if necessary, we can move the term containing $\epsilon M_n$ to the left and use $a+S_{n-1}\le(1+a/B_0)S_n$ to obtain $M_n\le K_p S_n\sqrt{r_n}$.
Also, \eqref{any:eq:total-sum-lower} gives $S_n\ge1/(8r_n)$ when $r_n\le1/8$.
Combining these estimates with \cref{any:lem:record-sum} yields
\[
    S_n^\nu\le K_p n r_n^{(1-\nu)/2}, \qquad 1\le K_p n r_n^{(1+\nu)/2}.
\]
Consequently, at every such record time,
\begin{equation}
    r_n\ge c_p n^{-\beta(p)}, \qquad \beta(p):=\frac{2p}{1+p}=\frac{2}{1+\nu}.
    \label{any:eq:fixed-p-rate}
\end{equation}

It remains to let $p$ approach $\psil$.
Write $\delta:=p-\psil$.
The parameter bound \eqref{eq:parameter-growth} and the constants in the proof of \cref{any:lem:record-sum} imply
\begin{equation}
\begin{aligned}
    \log(1/c_p)+\log(1/r_{0,p}) &=O\!\left(\delta^{-1}\log(1/\delta)\right),\\
    \log B_0&=O\!\left(\log(1/\delta)\right), \qquad \beta(p)=\pany+O(\delta).
\end{aligned}
\label{any:eq:parameter-growth}
\end{equation}
At each sufficiently large strict record time, choose
\[
    \delta=\sqrt{\frac{\log\log n}{\log n}}.
\]
The hard function may depend on $n$, since $\mathcal R_n(H_n)$ takes a separate supremum at each horizon.
If $M_n<B_0$, then $S_n\le nB_0$, and \eqref{any:eq:total-sum-lower} gives $r_n\ge n^{-1-o(1)}$.
If $r_n>r_{0,p}$, then \eqref{any:eq:parameter-growth} gives $r_n\ge n^{-o(1)}$.
Both cases are stronger than the claimed bound.
Otherwise, \eqref{any:eq:fixed-p-rate} and \eqref{any:eq:parameter-growth} give
\[
    r_n\ge n^{-\pany} \exp\!\left\{-C\left(\delta\log n+\delta^{-1}\log(1/\delta)\right)\right\} \ge n^{-\left(\pany+C'\sqrt{\log\log n/\log n}\right)}.
\]
There are infinitely many strict record times, which proves \cref{thm:anytime}.

\end{document}
